
\Data\ID{1003109301}
\Updated{2022-04-23 05:04:23}
\Level{A}
\SubArea{Binomial equations}
\LANGen{
\Question{
Consider the equation \( x^4 + 81 = 0 \). Find the sum of all its roots in the set of complex numbers.}
{\( 0 \)}{\( 3 \)}{\( -3 \)}{\( 81 \)}{\( - 81 \)}{}}
\LANGpl{
\Question{
Dane jest równanie \( x^4 + 81 = 0 \). Oblicz sumę wszystkich jego pierwiastków w zbiorze liczb zespolonych.}
{\( 0 \)}{\( 3 \)}{\( -3 \)}{\( 81 \)}{\( - 81 \)}{}}
\LANGcs{
\Question{
Součet kořenů binomické rovnice \( x^4 + 81 = 0 \) je:}
{\( 0 \)}{\( 3 \)}{\( -3 \)}{\( 81 \)}{\( - 81 \)}{}}
\LANGsk{
\Question{
Súčet koreňov binomickej rovnice \( x^4 + 81 = 0 \) je:}
{\( 0 \)}{\( 3 \)}{\( -3 \)}{\( 81 \)}{\( - 81 \)}{}}
\LANGes{
\Question{
Considera la ecuación \( x^4 + 81 = 0 \). Determina la suma de todas sus raíces en el conjunto de los números complejos.}
{\( 0 \)}{\( 3 \)}{\( -3 \)}{\( 81 \)}{\( - 81 \)}{}}
\End

\Data\ID{1003109302}
\Updated{2021-04-25 06:04:08}
\Level{A}
\SubArea{Binomial equations}
\LANGen{
\Question{
Consider the equation \( x^4 - 81 = 0 \). Find the product of all its roots in the set of complex numbers.}
{\( -81 \)}{\( 81 \)}{\( -3 \)}{\( 3 \)}{\( 0 \)}{}}
\LANGpl{
\Question{
Dane jest równanie  \( x^4 - 81 = 0 \). Oblicz iloczyn wszystkich jego pierwiastków w zbiorze liczb zespolonych.}
{\( -81 \)}{\( 81 \)}{\( -3 \)}{\( 3 \)}{\( 0 \)}{}}
\LANGcs{
\Question{
Součin kořenů binomické rovnice \( x^4 - 81 = 0 \) je:}
{\( -81 \)}{\( 81 \)}{\( -3 \)}{\( 3 \)}{\( 0 \)}{}}
\LANGsk{
\Question{
Súčin koreňov binomickej rovnice \( x^4 - 81 = 0 \) je:}
{\( -81 \)}{\( 81 \)}{\( -3 \)}{\( 3 \)}{\( 0 \)}{}}
\LANGes{
\Question{
Considera la ecuación \( x^4 - 81 = 0 \). Determina el producto de todas sus raíces en el conjunto de los números complejos.}
{\( -81 \)}{\( 81 \)}{\( -3 \)}{\( 3 \)}{\( 0 \)}{}}
\End

\Data\ID{1003109305}
\Updated{2022-04-23 05:04:23}
\Level{A}
\SubArea{Binomial equations}
\LANGen{
\Question{
Solve the following equation in the set of complex numbers. (Solve the equation by substitution.)
\[ (2x + 3)^4 - 256 = 0 \]}
{\( \left\{-\frac72;\frac12;-\frac32\pm2\mathrm{i} \right\} \)}{\( \left\{-\frac72;\frac12;\frac32\pm2\mathrm{i} \right\} \)}{\( \left\{\frac72;-\frac12;\frac32\pm2\mathrm{i} \right\} \)}{\( \left\{\frac72;-\frac12;-\frac32\pm2\mathrm{i} \right\} \)}{}{}}
\LANGpl{
\Question{
Rozwiąż  równanie w zbiorze liczb zespolonych. (Rozwiąż równanie przez podstawienie.)
\[ (2x + 3)^4 - 256 = 0 \]}
{\( \left\{-\frac72;\frac12;-\frac32\pm2\mathrm{i} \right\} \)}{\( \left\{-\frac72;\frac12;\frac32\pm2\mathrm{i} \right\} \)}{\( \left\{\frac72;-\frac12;\frac32\pm2\mathrm{i} \right\} \)}{\( \left\{\frac72;-\frac12;-\frac32\pm2\mathrm{i} \right\} \)}{}{}}
\LANGcs{
\Question{
Najděte množinu komplexních kořenů binomické rovnice \[ (2x + 3)^4 - 256 = 0. \]  (Řešte pomocí substituce.)}
{\( \left\{-\frac72;\frac12;-\frac32\pm2\mathrm{i} \right\} \)}{\( \left\{-\frac72;\frac12;\frac32\pm2\mathrm{i} \right\} \)}{\( \left\{\frac72;-\frac12;\frac32\pm2\mathrm{i} \right\} \)}{\( \left\{\frac72;-\frac12;-\frac32\pm2\mathrm{i} \right\} \)}{}{}}
\LANGsk{
\Question{
Nájdite množinu komplexných koreňov binomickej rovnice \[ (2x + 3)^4 - 256 = 0. \]  (Riešte pomocou substitúcie.)}
{\( \left\{-\frac72;\frac12;-\frac32\pm2\mathrm{i} \right\} \)}{\( \left\{-\frac72;\frac12;\frac32\pm2\mathrm{i} \right\} \)}{\( \left\{\frac72;-\frac12;\frac32\pm2\mathrm{i} \right\} \)}{\( \left\{\frac72;-\frac12;-\frac32\pm2\mathrm{i} \right\} \)}{}{}}
\LANGes{
\Question{
Resuelve la siguiente ecuación en el conjunto de los números complejos. (Resuelve la ecuación por el método de sustitución).
\[ (2x + 3)^4 - 256 = 0 \]}
{\( \left\{-\frac72;\frac12;-\frac32\pm2\mathrm{i} \right\} \)}{\( \left\{-\frac72;\frac12;\frac32\pm2\mathrm{i} \right\} \)}{\( \left\{\frac72;-\frac12;\frac32\pm2\mathrm{i} \right\} \)}{\( \left\{\frac72;-\frac12;-\frac32\pm2\mathrm{i} \right\} \)}{}{}}
\End

\Data\ID{1003118401}
\Updated{2021-04-25 06:04:59}
\Level{A}
\SubArea{Binomial equations}
\LANGen{
\Question{
Find the solution set of the equation \( x^3 - 8\mathrm{i} = 0 \) in the set of complex numbers.}
{\( \left\{\sqrt3+\mathrm{i}; -\sqrt3+\mathrm{i};-2\mathrm{i} \right\} \)}{\( \left\{ 2\mathrm{i}; -\sqrt3-\mathrm{i}; \sqrt3-\mathrm{i} \right\} \)}{\( \left\{\frac{\sqrt3}2+\frac12\mathrm{i}; -\frac{\sqrt3}2+\frac12\mathrm{i};-\mathrm{i} \right\} \)}{\( \left\{\mathrm{i};-\frac{\sqrt3}2-\frac12\mathrm{i}; \frac{\sqrt3}2-\frac12\mathrm{i} \right\} \)}{}{}}
\LANGpl{
\Question{
Wskaż zbiór rozwiązań równania  \( x^3 - 8\mathrm{i} = 0 \) w zbiorze liczb zespolonych.}
{\( \left\{\sqrt3+\mathrm{i}; -\sqrt3+\mathrm{i};-2\mathrm{i} \right\} \)}{\( \left\{ 2\mathrm{i}; -\sqrt3-\mathrm{i}; \sqrt3-\mathrm{i} \right\} \)}{\( \left\{\frac{\sqrt3}2+\frac12\mathrm{i}; -\frac{\sqrt3}2+\frac12\mathrm{i};-\mathrm{i} \right\} \)}{\( \left\{\mathrm{i};-\frac{\sqrt3}2-\frac12\mathrm{i}; \frac{\sqrt3}2-\frac12\mathrm{i} \right\} \)}{}{}}
\LANGcs{
\Question{
Množina komplexních kořenů binomické rovnice \( x^3 - 8\mathrm{i} = 0 \) je:}
{\( \left\{\sqrt3+\mathrm{i}; -\sqrt3+\mathrm{i};-2\mathrm{i} \right\} \)}{\( \left\{ 2\mathrm{i}; -\sqrt3-\mathrm{i}; \sqrt3-\mathrm{i} \right\} \)}{\( \left\{\frac{\sqrt3}2+\frac12\mathrm{i}; -\frac{\sqrt3}2+\frac12\mathrm{i};-\mathrm{i} \right\} \)}{\( \left\{\mathrm{i};-\frac{\sqrt3}2-\frac12\mathrm{i}; \frac{\sqrt3}2-\frac12\mathrm{i} \right\} \)}{}{}}
\LANGsk{
\Question{
Množina komplexných koreňov binomickej rovnice \( x^3 - 8\mathrm{i} = 0 \) je:}
{\( \left\{\sqrt3+\mathrm{i}; -\sqrt3+\mathrm{i};-2\mathrm{i} \right\} \)}{\( \left\{ 2\mathrm{i}; -\sqrt3-\mathrm{i}; \sqrt3-\mathrm{i} \right\} \)}{\( \left\{\frac{\sqrt3}2+\frac12\mathrm{i}; -\frac{\sqrt3}2+\frac12\mathrm{i};-\mathrm{i} \right\} \)}{\( \left\{\mathrm{i};-\frac{\sqrt3}2-\frac12\mathrm{i}; \frac{\sqrt3}2-\frac12\mathrm{i} \right\} \)}{}{}}
\LANGes{
\Question{
Determina el conjunto de soluciones de la ecuación \( x^3 - 8\mathrm{i} = 0 \) en el conjunto de los números complejos.}
{\( \left\{\sqrt3+\mathrm{i}; -\sqrt3+\mathrm{i};-2\mathrm{i} \right\} \)}{\( \left\{ 2\mathrm{i}; -\sqrt3-\mathrm{i}; \sqrt3-\mathrm{i} \right\} \)}{\( \left\{\frac{\sqrt3}2+\frac12\mathrm{i}; -\frac{\sqrt3}2+\frac12\mathrm{i};-\mathrm{i} \right\} \)}{\( \left\{\mathrm{i};-\frac{\sqrt3}2-\frac12\mathrm{i}; \frac{\sqrt3}2-\frac12\mathrm{i} \right\} \)}{}{}}
\End

\Data\ID{1003118402}
\Updated{2020-07-19 02:07:52}
\Level{A}
\SubArea{Binomial equations}
\LANGen{
\Question{
Which of the given complex numbers is not a root of the equation \( x^6  +  8\mathrm{i}  =  0 \)?}
{\( 1-\mathrm{i} \)}{\( 1+\mathrm{i} \)}{\( -1-\mathrm{i} \)}{\( \sqrt2\left(\cos\frac{7\pi}{12}+\mathrm{i}\cdot\sin\frac{7\pi}{12}\right) \)}{\( \sqrt2\left(\cos\frac{23\pi}{12}+\mathrm{i}\cdot\sin\frac{23\pi}{12}\right) \)}{}}
\LANGpl{
\Question{
Która z podanych liczb zespolonych nie jest pierwiastkiem równania \( x^6  +  8\mathrm{i}  =  0 \)?}
{\( 1-\mathrm{i} \)}{\( 1+\mathrm{i} \)}{\( -1-\mathrm{i} \)}{\( \sqrt2\left(\cos\frac{7\pi}{12}+\mathrm{i}\cdot\sin\frac{7\pi}{12}\right) \)}{\( \sqrt2\left(\cos\frac{23\pi}{12}+\mathrm{i}\cdot\sin\frac{23\pi}{12}\right) \)}{}}
\LANGcs{
\Question{
Kořenem binomické rovnice \( x^6  +  8\mathrm{i}  =  0 \) není:}
{\( 1-\mathrm{i} \)}{\( 1+\mathrm{i} \)}{\( -1-\mathrm{i} \)}{\( \sqrt2\left(\cos\frac{7\pi}{12}+\mathrm{i}\cdot\sin\frac{7\pi}{12}\right) \)}{\( \sqrt2\left(\cos\frac{23\pi}{12}+\mathrm{i}\cdot\sin\frac{23\pi}{12}\right) \)}{}}
\LANGsk{
\Question{
Koreňom binomickej rovnice \( x^6  +  8\mathrm{i}  =  0 \) nie je:}
{\( 1-\mathrm{i} \)}{\( 1+\mathrm{i} \)}{\( -1-\mathrm{i} \)}{\( \sqrt2\left(\cos\frac{7\pi}{12}+\mathrm{i}\cdot\sin\frac{7\pi}{12}\right) \)}{\( \sqrt2\left(\cos\frac{23\pi}{12}+\mathrm{i}\cdot\sin\frac{23\pi}{12}\right) \)}{}}
\LANGes{
\Question{
¿Cuál de los números complejos dados no es una raíz de la ecuación \( x^6  +  8\mathrm{i}  =  0 \)?}
{\( 1-\mathrm{i} \)}{\( 1+\mathrm{i} \)}{\( -1-\mathrm{i} \)}{\( \sqrt2\left(\cos\frac{7\pi}{12}+\mathrm{i}\cdot\sin\frac{7\pi}{12}\right) \)}{\( \sqrt2\left(\cos\frac{23\pi}{12}+\mathrm{i}\cdot\sin\frac{23\pi}{12}\right) \)}{}}
\End

\Data\ID{1003118405}
\Updated{2022-04-23 05:04:23}
\Level{A}
\SubArea{Binomial equations}
\LANGen{
\Question{
All the solutions of the equation \( x^6-4\sqrt3+4\mathrm{i} = 0 \) can be  displayed as points in the  rectangular coordinate system. What is the distance of the two most distant points?}
{\( 2\sqrt2 \)}{\( \sqrt2 \)}{\( 2\sqrt[3]4 \)}{\( \sqrt[3]4 \)}{\( 2\sqrt3 \)}{\( \sqrt3 \)}}
\LANGpl{
\Question{
Wszystkie rozwiązania równania  \( x^6-4\sqrt3+4\mathrm{i} = 0 \) mogą być zaznaczone  jako punkty w  układzie współrzędnych. Jaka jest odległość dwóch najbardziej odległych punktów?}
{\( 2\sqrt2 \)}{\( \sqrt2 \)}{\( 2\sqrt[3]4 \)}{\( \sqrt[3]4 \)}{\( 2\sqrt3 \)}{\( \sqrt3 \)}}
\LANGcs{
\Question{
Vzdálenost dvou nejvzdálenějších obrazů kořenů binomické rovnice  \( x^6-4\sqrt3+4\mathrm{i} = 0 \), zobrazených v Gaussově rovině, je:}
{\( 2\sqrt2 \)}{\( \sqrt2 \)}{\( 2\sqrt[3]4 \)}{\( \sqrt[3]4 \)}{\( 2\sqrt3 \)}{\( \sqrt3 \)}}
\LANGsk{
\Question{
Vzdialenosť dvoch najvzdialenejších obrazov koreňov binomickej rovnice  \( x^6-4\sqrt3+4\mathrm{i} = 0 \), zobrazených v Gaussovej rovine, je:}
{\( 2\sqrt2 \)}{\( \sqrt2 \)}{\( 2\sqrt[3]4 \)}{\( \sqrt[3]4 \)}{\( 2\sqrt3 \)}{\( \sqrt3 \)}}
\LANGes{
\Question{
Todas las soluciones de la ecuación \( x^6-4\sqrt3+4\mathrm{i} = 0 \) se pueden mostrar como puntos en el sistema de coordenadas rectangular. ¿Cuál es la distancia de los dos puntos más distantes?}
{\( 2\sqrt2 \)}{\( \sqrt2 \)}{\( 2\sqrt[3]4 \)}{\( \sqrt[3]4 \)}{\( 2\sqrt3 \)}{\( \sqrt3 \)}}
\End

\Data\ID{1003118406}
\Updated{2021-04-25 06:04:00}
\Level{A}
\SubArea{Binomial equations}
\LANGen{
\Question{
All the solutions of the equation \( x^4+1+\sqrt3\mathrm{i} = 0 \) are complex numbers, whose arguments are from the interval \( [0; 2\pi) \). Find the sum of the arguments of all the solutions of the equation.}
{\( \frac{13}3\pi \)}{\( 4\pi \)}{\( \frac{25}6\pi \)}{\( \frac92\pi \)}{}{}}
\LANGpl{
\Question{
Wszystkie rozwiązania równania \( x^4+1+\sqrt3\mathrm{i} = 0 \) są liczbami zespolonymi, których argumenty należą  do  przedziału \( \langle0; 2\pi) \). Znajdź sumę argumentów wszystkich rozwiązań równania.}
{\( \frac{13}3\pi \)}{\( 4\pi \)}{\( \frac{25}6\pi \)}{\( \frac92\pi \)}{}{}}
\LANGcs{
\Question{
Všechny kořeny rovnice  \( x^4+1+\sqrt3\mathrm{i} = 0 \) jsou komplexní čísla s argumenty z intervalu \( \langle0; 2\pi) \). Určete součet argumentů všech kořenů rovnice.}
{\( \frac{13}3\pi \)}{\( 4\pi \)}{\( \frac{25}6\pi \)}{\( \frac92\pi \)}{}{}}
\LANGsk{
\Question{
Všetky korene rovnice  \( x^4+1+\sqrt3\mathrm{i} = 0 \) sú komplexné čísla s argumentami z intervalu \( \langle0; 2\pi) \). Určte súčet argumentov všetkých koreňov rovnice.}
{\( \frac{13}3\pi \)}{\( 4\pi \)}{\( \frac{25}6\pi \)}{\( \frac92\pi \)}{}{}}
\LANGes{
\Question{
Todas las soluciones de la ecuación \( x^4+1+\sqrt3\mathrm{i} = 0 \) son números complejos, cuyos argumentos pertenecen al intervalo \( [0; 2\pi) \). Calcula la suma de los argumentos de todas las soluciones de la ecuación.}
{\( \frac{13}3\pi \)}{\( 4\pi \)}{\( \frac{25}6\pi \)}{\( \frac92\pi \)}{}{}}
\End

\Data\ID{1103109303}
\Updated{2020-07-20 02:07:26}
\Level{A}
\SubArea{Binomial equations}
\IMG{1103109303.pdf}
\LANGen{
\Question{
Consider an equation \( x^n+b=0 \), where \( n \) is a positive integer and \( b \) is a real number. The points that correspond to the roots of the equation are marked in the figure as black points. Find the equation.}
{\( x^8 - 256 = 0 \)}{\( x^8 + 256 = 0 \)}{\( x^4 + 16 = 0 \)}{\( x^4 - 16 = 0 \)}{\( x^6 - 64 = 0 \)}{\( x^6 + 64 = 0 \)}}
\LANGpl{
\Question{
Dane jest równanie \( x^n+b=0 \), gdzie \( n \) jest  dodatnią liczbą całkowitą, a  \( b \) jest liczbą rzeczywistą.  Punkty odpowiadające pierwiastkom równania są oznaczone na rysunku kolorem czarnym. Wskaż  równanie.}
{\( x^8 - 256 = 0 \)}{\( x^8 + 256 = 0 \)}{\( x^4 + 16 = 0 \)}{\( x^4 - 16 = 0 \)}{\( x^6 - 64 = 0 \)}{\( x^6 + 64 = 0 \)}}
\LANGcs{
\Question{
Na obrázku jsou černými body zobrazeny kořeny binomické rovnice \( x^n+b=0 \), kde \( n \) je přirozené číslo a \( b \) je reálné číslo. Určete tuto rovnici.}
{\( x^8 - 256 = 0 \)}{\( x^8 + 256 = 0 \)}{\( x^4 + 16 = 0 \)}{\( x^4 - 16 = 0 \)}{\( x^6 - 64 = 0 \)}{\( x^6 + 64 = 0 \)}}
\LANGsk{
\Question{
Na obrázku sú čiernymi bodmi zobrazené korene binomickej rovnice \( x^n+b=0 \), kde \( n \) je prirodzené číslo a \( b \) je reálne číslo. Určte túto rovnicu.}
{\( x^8 - 256 = 0 \)}{\( x^8 + 256 = 0 \)}{\( x^4 + 16 = 0 \)}{\( x^4 - 16 = 0 \)}{\( x^6 - 64 = 0 \)}{\( x^6 + 64 = 0 \)}}
\LANGes{
\Question{
Considera la ecuación \( x^n+b=0 \), donde \( n \) es un número natural y \( b \) es un número real. Los puntos que corresponden a las raíces de la ecuación están marcados en la figura como puntos negros. Determina la ecuación.}
{\( x^8 - 256 = 0 \)}{\( x^8 + 256 = 0 \)}{\( x^4 + 16 = 0 \)}{\( x^4 - 16 = 0 \)}{\( x^6 - 64 = 0 \)}{\( x^6 + 64 = 0 \)}}
\End

\Data\ID{1103109304}
\Updated{2020-07-20 02:07:43}
\Level{A}
\SubArea{Binomial equations}
\LANGen{
\Question{
Choose the figure in which the points marked in black correspond to the roots of the equation  \( x^5 + 1 = 0 \).}
{}{}{}{}{}{}}
\LANGpl{
\Question{
Wskaż figurę, na której punkty zaznaczone na czarno odpowiadają pierwiastkom  równania \( x^5 + 1 = 0 \).}
{}{}{}{}{}{}}
\LANGcs{
\Question{
Na kterém obrázku jsou černými body zobrazeny kořeny rovnice  \( x^5 + 1 = 0 \)?}
{}{}{}{}{}{}}
\LANGsk{
\Question{
Na ktorom obrázku sú čiernymi bodmi zobrazené korene binomickej rovnice  \( x^5 + 1 = 0 \)?}
{}{}{}{}{}{}}
\LANGes{
\Question{
Elige la figura en la que los puntos marcados en negro corresponden a las raíces de la ecuación \( x^5 + 1 = 0 \).}
{}{}{}{}{}{}}

\IMGanswer{1103109304a.pdf}
\IMGanswer{1103109304b.pdf}
\IMGanswer{1103109304c.pdf}
\IMGanswer{1103109304d.pdf}\End

\Data\ID{1103118403}
\Updated{2020-07-19 02:07:54}
\Level{A}
\SubArea{Binomial equations}
\LANGen{
\Question{
Consider the equation  \( x^4+2-2\sqrt3\mathrm{i} = 0 \).  In one of the following pictures the points depicted in black correspond to the roots of the equation. Find the figure.}
{}{}{}{}{}{}}
\LANGpl{
\Question{
Dane jest równanie  \( x^4+2-2\sqrt3\mathrm{i} = 0 \).  Na jednym z poniższych rysunków punkty oznaczone kolorem czarnym odpowiadają pierwiastkom danego równania. Wskaż figurę.}
{}{}{}{}{}{}}
\LANGcs{
\Question{
Vyberte obrázek, na kterém jsou černými body zobrazeny kořeny binomické rovnice:  \( x^4+2-2\sqrt3\mathrm{i} = 0 \).}
{}{}{}{}{}{}}
\LANGsk{
\Question{
Vyberte obrázok, na ktorom sú čiernymi bodmi zobrazené korene binomickej rovnice:  \( x^4+2-2\sqrt3\mathrm{i} = 0 \).}
{}{}{}{}{}{}}
\LANGes{
\Question{
Considera la ecuación \( x^4+2-2\sqrt3\mathrm{i} = 0 \).  En una de las siguientes figuras, los puntos representados en negro corresponden a las raíces de la ecuación. Determina la figura.}
{}{}{}{}{}{}}

\IMGanswer{1103118403a_0.pdf}
\IMGanswer{1103118403b_0.pdf}
\IMGanswer{1103118403c_0.pdf}
\IMGanswer{1103118403d_0.pdf}\End

\Data\ID{1103118404}
\Updated{2020-07-19 02:07:33}
\Level{A}
\SubArea{Binomial equations}
\IMG{1103118404_0.pdf}
\LANGen{
\Question{
Consider an equation \( x^n+b=0 \), where \( n \) is a positive integer and \( b \) is a complex number. In the picture the points that correspond to the roots of the equation are depicted in black. Find the equation.}
{\( x^3 + 4\sqrt2 - 4\sqrt2\mathrm{i} = 0 \)}{\( x^3 + 4\sqrt2 +4\sqrt2\mathrm{i} = 0 \)}{\( x^3 - 4\sqrt2 - 4\sqrt2\mathrm{i} = 0 \)}{\( x^3 - 4\sqrt2 +4\sqrt2\mathrm{i} = 0 \)}{}{}}
\LANGpl{
\Question{
Dane jest równanie \( x^n+b=0 \), gdzie \( n \) jest dodatnią liczbą naturalna, a  \( b \) jest liczbą zespoloną. Na rysunku punkty oznaczone kolorem czarnym odpowiadają pierwiastkom równania. Wskaż równanie.}
{\( x^3 + 4\sqrt2 - 4\sqrt2\mathrm{i} = 0 \)}{\( x^3 + 4\sqrt2 +4\sqrt2\mathrm{i} = 0 \)}{\( x^3 - 4\sqrt2 - 4\sqrt2\mathrm{i} = 0 \)}{\( x^3 - 4\sqrt2 +4\sqrt2\mathrm{i} = 0 \)}{}{}}
\LANGcs{
\Question{
Uvažujte rovnici $x^n+b=0$, kde $n\in\mathbb{N}^+$ a $b$ je komplexní číslo. Na obrázku jsou černými body zobrazeny kořeny binomické rovnice:}
{\( x^3 + 4\sqrt2 - 4\sqrt2\mathrm{i} = 0 \)}{\( x^3 + 4\sqrt2 +4\sqrt2\mathrm{i} = 0 \)}{\( x^3 - 4\sqrt2 - 4\sqrt2\mathrm{i} = 0 \)}{\( x^3 - 4\sqrt2 +4\sqrt2\mathrm{i} = 0 \)}{}{}}
\LANGsk{
\Question{
Uvažujte rovnicu $x^n+b=0$, kde $n\in\mathbb{N}^+$ a $b$ je komplexné číslo. Na obrázku sú čiernymi bodmi zobrazené korene binomickej rovnice:}
{\( x^3 + 4\sqrt2 - 4\sqrt2\mathrm{i} = 0 \)}{\( x^3 + 4\sqrt2 +4\sqrt2\mathrm{i} = 0 \)}{\( x^3 - 4\sqrt2 - 4\sqrt2\mathrm{i} = 0 \)}{\( x^3 - 4\sqrt2 +4\sqrt2\mathrm{i} = 0 \)}{}{}}
\LANGes{
\Question{
Considera la ecuación \( x^n+b=0 \), donde \( n \) es un número natural y \( b \) es un número complejo.  En la imagen, los puntos que corresponden a las raíces de la ecuación se representan en negro. Determina la ecuación.}
{\( x^3 + 4\sqrt2 - 4\sqrt2\mathrm{i} = 0 \)}{\( x^3 + 4\sqrt2 +4\sqrt2\mathrm{i} = 0 \)}{\( x^3 - 4\sqrt2 - 4\sqrt2\mathrm{i} = 0 \)}{\( x^3 - 4\sqrt2 +4\sqrt2\mathrm{i} = 0 \)}{}{}}
\End

\Data\ID{2000002601}
\Updated{2021-12-01 04:12:54}
\Level{A}
\SubArea{Binomial equations}
\LANGen{
\Question{
Consider the equation \(x^4+16=0\). Which of the given numbers is the solution of this equation?}
{\( 2(\cos{\frac{\pi}{4}}+ i \sin{\frac{\pi}{4}}) \)}{\( 2i\)}{\( -2i\)}{\( 2(\cos{{\pi}}+ i \sin{{\pi}}) \)}{}{}}
\LANGpl{
\Question{
Wskaż rozwiązanie równania \(x^4+16=0\).}
{\( 2(\cos{\frac{\pi}{4}}+ i \sin{\frac{\pi}{4}}) \)}{\( 2i\)}{\( -2i\)}{\( 2(\cos{{\pi}}+ i \sin{{\pi}}) \)}{}{}}
\LANGcs{
\Question{
Mějme rovnici \(x^4+16=0\). Které z daných čísel je řešením této rovnice?}
{\( 2(\cos{\frac{\pi}{4}}+ i \sin{\frac{\pi}{4}}) \)}{\( 2i\)}{\( -2i\)}{\( 2(\cos{{\pi}}+ i \sin{{\pi}}) \)}{}{}}
\LANGsk{
\Question{
Ktoré z daných čísel je riešením rovnice \(x^4+16=0\)?}
{\( 2(\cos{\frac{\pi}{4}}+ i \sin{\frac{\pi}{4}}) \)}{\( 2i\)}{\( -2i\)}{\( 2(\cos{{\pi}}+ i \sin{{\pi}}) \)}{}{}}
\LANGes{
\Question{
Considera la ecuación \(x^4+16=0\). ¿Cuál de los números dados es la solución de esta ecuación?}
{\( 2(\cos{\frac{\pi}{4}}+ i \sin{\frac{\pi}{4}}) \)}{\( 2i\)}{\( -2i\)}{\( 2(\cos{{\pi}}+ i \sin{{\pi}}) \)}{}{}}
\End

\Data\ID{2000002602}
\Updated{2021-12-26 08:12:43}
\Level{A}
\SubArea{Binomial equations}
\LANGen{
\Question{
Consider the equation  \(x^4 =1\), where \(x\) is a complex variable. Which of the following statements is true?}
{The equation has four different complex roots.}{The equation has no real root.}{The equation has two double roots:  \(x_{1,2}=1\) and \(x_{3,4}=-1\).}{The equation has the root \(x=1+i\).}{}{}}
\LANGpl{
\Question{
Dane jest równanie   \(x^4 =1\),  gdzie \(x\) jest zmienną zespoloną. Wybierz  poprawne stwierdzenie.}
{Równanie ma cztery różne rozwiązania.}{Równanie nie ma rozwiązania.}{Równanie ma dwa rozwiązania:  \(x_{1,2}=1\)  i  \(x_{3,4}=-1\).}{Rozwiązaniem równania jest \(x=1+i\).}{}{}}
\LANGcs{
\Question{
Mějme rovnici \(x^4 =1\), kde \(x\) je komplexní proměnná. Které z následujících tvrzení je pravdivé?}
{Rovnice má čtyři různé komplexní kořeny.}{Rovnice nemá reálný kořen.}{Rovnice má dva dvojnásobné kořeny:  \(x_{1,2}=1\) a \(x_{3,4}=-1\).}{Rovnice má kořen \(x=1+i\).}{}{}}
\LANGsk{
\Question{
Daná je rovnica \(x^4 =1\), kde \(x\)  je komplexná premenná. Ktoré z nasledujúcich tvrdení o riešení rovnice je pravdivé?}
{Rovnica má štyri rôzne komplexné korene.}{Rovnica nemá reálny koreň.}{Rovnica má dva dvojnásobné korene: \(x_{1,2}=1\) a \(x_{3,4}=-1\).}{Rovnica má koreň \(x=1+i\).}{}{}}
\LANGes{
\Question{
Considera la ecuación  \(x^4 =1\), donde \(x\) es una variable compleja. ¿Cuál de las siguientes proposiciones es verdadera?}
{La ecuación tiene cuatro raíces complejas diferentes.}{La ecuación no tiene ninguna raíz real.}{La ecuación tiene dos raíces dobles:  \(x_{1,2}=1\) y \(x_{3,4}=-1\).}{La ecuación tiene como raíz \(x=1+i\).}{}{}}
\End

\Data\ID{2000002603}
\Updated{2021-12-01 04:12:07}
\Level{A}
\SubArea{Binomial equations}
\LANGen{
\Question{
One of the roots of the equation \(x^3-8=0\) is \(x_1 = -1-i\sqrt{3}\). Find the sum of all its roots.}
{\( 0\)}{\( -8\)}{\( -2i\sqrt{3} \)}{\(-4\)}{}{}}
\LANGpl{
\Question{
Jednym z pierwiastków równania \(x^3-8=0\)  jest \(x_1 = -1-i\sqrt{3}\). Znajdź sumę wszystkich jego pierwiastków.}
{\( 0\)}{\( -8\)}{\( -2i\sqrt{3} \)}{\(-4\)}{}{}}
\LANGcs{
\Question{
Jedním z kořenů rovnice \(x^3-8=0\) je \(x_1 = -1-i\sqrt{3}\). Určete součet všech kořenů rovnice.}
{\( 0\)}{\( -8\)}{\( -2i\sqrt{3} \)}{\(-4\)}{}{}}
\LANGsk{
\Question{
Jeden z koreňov rovnice \(x^3-8=0\) je \(x_1 = -1-i\sqrt{3}\). 
Určte súčet všetkých koreňov tejto rovnice.}
{\( 0\)}{\( -8\)}{\( -2i\sqrt{3} \)}{\(-4\)}{}{}}
\LANGes{
\Question{
Una de las raíces de la ecuación \(x^3-8=0\) es \(x_1 = -1-i\sqrt{3}\). Calcula la suma de todas sus raíces.}
{\( 0\)}{\( -8\)}{\( -2i\sqrt{3} \)}{\(-4\)}{}{}}
\End

\Data\ID{2000002604}
\Updated{2021-12-01 04:12:11}
\Level{A}
\SubArea{Binomial equations}
\LANGen{
\Question{
Find the solution set of the equation \(x^4+81=0\) if you know that one of its roots is \(\frac{3}{\sqrt{2}}(1+i)\).}
{\( \left\{ \frac{3}{\sqrt{2}}(1+i); -\frac{3}{\sqrt{2}}(1+i); \frac{3}{\sqrt{2}}(1-i);-\frac{3}{\sqrt{2}}(1-i) \right\} \)}{\( \left\{ \frac{3}{\sqrt{2}}(1+i); -\frac{3}{\sqrt{2}}(1+i);3;-3 \right\} \)}{\( \left\{ \frac{3}{\sqrt{2}}(1+i); \frac{3}{\sqrt{2}}(1-i);3i;-3i \right\} \)}{\( \left\{\frac{3}{\sqrt{2}}(1+i);\frac{3}{\sqrt{2}}(1-i) \right\}\)}{}{}}
\LANGpl{
\Question{
Wskaż zbiór rozwiązań równania \(x^4+81=0\) wiedząc, że jednym z jego pierwiastków  jest \(\frac{3}{\sqrt{2}}(1+i)\).}
{\( \left\{ \frac{3}{\sqrt{2}}(1+i); -\frac{3}{\sqrt{2}}(1+i); \frac{3}{\sqrt{2}}(1-i);-\frac{3}{\sqrt{2}}(1-i) \right\} \)}{\( \left\{ \frac{3}{\sqrt{2}}(1+i); -\frac{3}{\sqrt{2}}(1+i);3;-3 \right\} \)}{\( \left\{ \frac{3}{\sqrt{2}}(1+i); \frac{3}{\sqrt{2}}(1-i);3i;-3i \right\} \)}{\( \left\{\frac{3}{\sqrt{2}}(1+i);\frac{3}{\sqrt{2}}(1-i) \right\}\)}{}{}}
\LANGcs{
\Question{
Určete množinu řešení rovnice \(x^4+81=0\), když víte, že jedním z kořenů je \(\frac{3}{\sqrt{2}}(1+i)\).}
{\( \left\{ \frac{3}{\sqrt{2}}(1+i); -\frac{3}{\sqrt{2}}(1+i); \frac{3}{\sqrt{2}}(1-i);-\frac{3}{\sqrt{2}}(1-i) \right\} \)}{\( \left\{ \frac{3}{\sqrt{2}}(1+i); -\frac{3}{\sqrt{2}}(1+i);3;-3 \right\} \)}{\( \left\{ \frac{3}{\sqrt{2}}(1+i); \frac{3}{\sqrt{2}}(1-i);3i;-3i \right\} \)}{\( \left\{\frac{3}{\sqrt{2}}(1+i);\frac{3}{\sqrt{2}}(1-i) \right\}\)}{}{}}
\LANGsk{
\Question{
Určte množinu všetkých riešení rovnice \(x^4+81=0\), ak poznáte jeden z koreňov \(\frac{3}{\sqrt{2}}(1+i)\).}
{\( \left\{ \frac{3}{\sqrt{2}}(1+i); -\frac{3}{\sqrt{2}}(1+i); \frac{3}{\sqrt{2}}(1-i);-\frac{3}{\sqrt{2}}(1-i) \right\} \)}{\( \left\{ \frac{3}{\sqrt{2}}(1+i); -\frac{3}{\sqrt{2}}(1+i);3;-3 \right\} \)}{\( \left\{ \frac{3}{\sqrt{2}}(1+i); \frac{3}{\sqrt{2}}(1-i);3i;-3i \right\} \)}{\( \left\{\frac{3}{\sqrt{2}}(1+i);\frac{3}{\sqrt{2}}(1-i) \right\}\)}{}{}}
\LANGes{
\Question{
Determina el conjunto  solución de la ecuación \(x^4+81=0\) si sabes que una de sus raíces es\(\frac{3}{\sqrt{2}}(1+i)\).}
{\( \left\{ \frac{3}{\sqrt{2}}(1+i); -\frac{3}{\sqrt{2}}(1+i); \frac{3}{\sqrt{2}}(1-i);-\frac{3}{\sqrt{2}}(1-i) \right\} \)}{\( \left\{ \frac{3}{\sqrt{2}}(1+i); -\frac{3}{\sqrt{2}}(1+i);3;-3 \right\} \)}{\( \left\{ \frac{3}{\sqrt{2}}(1+i); \frac{3}{\sqrt{2}}(1-i);3i;-3i \right\} \)}{\( \left\{\frac{3}{\sqrt{2}}(1+i);\frac{3}{\sqrt{2}}(1-i) \right\}\)}{}{}}
\End

\Data\ID{2000002605}
\Updated{2021-12-01 04:12:13}
\Level{A}
\SubArea{Binomial equations}
\LANGen{
\Question{
How many solutions does the equation \(2x^4=32\) have in the set of complex numbers?}
{four}{one}{two}{eight}{}{}}
\LANGpl{
\Question{
Ile rozwiązań ma równanie \(2x^4=32\) w zbiorze liczb zespolonych?}
{cztery}{jedno}{dwa}{osiem}{}{}}
\LANGcs{
\Question{
Kolik řešení má rovnice \(2x^4=32\) na množině komplexních čísel?}
{čtyři}{jedno}{dvě}{osm}{}{}}
\LANGsk{
\Question{
Koľko riešení má rovnica \(2x^4=32\) v množine komplexných čísel?}
{štyri}{jedno}{dve}{osem}{}{}}
\LANGes{
\Question{
¿Cuántas soluciones tiene la ecuación \(2x^4=32\) en el conjunto de números complejos?}
{cuatro}{uno}{dos}{ocho}{}{}}
\End

\Data\ID{2000002606}
\Updated{2021-12-01 04:12:19}
\Level{A}
\SubArea{Binomial equations}
\LANGen{
\Question{
Imagine all the solutions of the equation \(x^6 -64 =0\) shown as points in the complex plane. Find the false statement.}
{Two points lie on the imaginary axis.}{The values of the arguments of any two solutions differ by an integer multiple of  \(\frac{\pi}{3}\).}{All solutions of the equation lie on a circle centered at the origin with a radius of \(2\).}{Two points lie on the real axis.}{}{}}
\LANGpl{
\Question{
Rozważ wszystkie rozwiązania równania \(x^6 -64 =0\)  przedstawione jako punkty na płaszczyźnie zespolonej. Wskaż fałszywe stwierdzenie.}
{Na osi urojonej leżą dwa punkty.}{Dwa rozwiązania różnią się całkowitą wielkokrotnością \(\frac{\pi}{3}\).}{Wszystkie rozwiązania równania leżą na okręgu o promieniu \(2\).}{Na osi rzeczywistej leżą dwa punkty.}{}{}}
\LANGcs{
\Question{
Řešení rovnice \(x^6 -64 =0\) jsou zobrazena jako body komplexní roviny. Vyberte nepravdivý výrok.}
{Dva body leží na imaginární ose.}{Hodnoty argumentů každých dvou řešení se liší o celočíselný násobek \(\frac{\pi}{3}\).}{Všechna řešení rovnice leží na kružnici se středem v počátku a poloměrem \(2\).}{Dva body leží na reálné ose.}{}{}}
\LANGsk{
\Question{
Všetky riešenia rovnice \(x^6 -64 =0\) sú zobrazené ako body komplexnej roviny. Vyberte nepravdivý výrok.}
{Dva body ležia na imaginárnej osi.}{Hodnoty argumentov každých dvoch riešení sa líšia o celočíselný násobok \(\frac{\pi}{3}\).}{Všetky riešenia rovnice ležia na kružnici so stredom v počiatku súradného systému s polomerom \(2\).}{Dva body ležia na reálnej osi.}{}{}}
\LANGes{
\Question{
Imagina que todas las soluciones de la ecuación \(x^6 -64 =0\) se muestran como puntos en el plano complejo. Determina la proposición falsa.}
{Dos puntos se encuentran en el eje imaginario.}{Los valores de los argumentos de dos soluciones cualesquiera  difieren en un múltiplo entero \(\frac{\pi}{3}\).}{Todas las soluciones de la ecuación se encuentran en un círculo centrado en el origen con un radio de \(2\).}{Dos puntos se encuentran en el eje real.}{}{}}
\End

\Data\ID{2000002607}
\Updated{2021-12-01 04:12:24}
\Level{A}
\SubArea{Binomial equations}
\LANGen{
\Question{
Which of the following equations has not the solution \(x=i\)?}
{\( x^6 -1 =0\)}{\( x^6 +1 =0\)}{\( x^3 +i =0\)}{\( x^5 -i=0\)}{}{}}
\LANGpl{
\Question{
Które z podanych równań nie ma rozwiązania \(x=i\)?}
{\( x^6 -1 =0\)}{\( x^6 +1 =0\)}{\( x^3 +i =0\)}{\( x^5 -i=0\)}{}{}}
\LANGcs{
\Question{
Která z následujících rovnic nemá řešení \(x=i\)?}
{\( x^6 -1 =0\)}{\( x^6 +1 =0\)}{\( x^3 +i =0\)}{\( x^5 -i=0\)}{}{}}
\LANGsk{
\Question{
Ktorá z následujúcich rovníc nemá riešenie \(x=i\)?}
{\( x^6 -1 =0\)}{\( x^6 +1 =0\)}{\( x^3 +i =0\)}{\( x^5 -i=0\)}{}{}}
\LANGes{
\Question{
¿Cuál de las siguientes ecuaciones no tiene la solución \(x=i\)?}
{\( x^6 -1 =0\)}{\( x^6 +1 =0\)}{\( x^3 +i =0\)}{\( x^5 -i=0\)}{}{}}
\End

\Data\ID{2000002608}
\Updated{2021-12-01 04:12:29}
\Level{A}
\SubArea{Binomial equations}
\LANGen{
\Question{
Find the right formula for solving the equation \(x^5 +32=0\).}
{\( x_k = \sqrt[5]{|-32|}( \cos\frac{\pi +2k\pi}{5}+ i\sin \frac{\pi +2k\pi}{5})\), \(k=0,1,2,3,4\)}{\( x_k = \sqrt[5]{-32}( \cos\frac{\pi +2k\pi}{5}+ i\sin \frac{\pi +2k\pi}{5})\), \(k=0,1,2,3,4\)}{\( x_k = \sqrt[5]{|-32|}( \cos \frac{\pi +k\pi}{5}+ i\sin \frac{\pi +k\pi}{5})\), \(k=0,1,2,3,4\)}{\( x_k = \sqrt[5]{|-32|}( \cos \frac{\pi +2k\pi}{5}+ \sin \frac{\pi +2k\pi}{5})\), \(k=0,1,2,3,4\)}{}{}}
\LANGpl{
\Question{
Wskaż odpowiedni wzór na rozwiązanie równania \(x^5 +32=0\).}
{\( x_k = \sqrt[5]{|-32|}( \cos\frac{\pi +2k\pi}{5}+ i\sin \frac{\pi +2k\pi}{5})\), \(k=0,1,2,3,4\)}{\( x_k = \sqrt[5]{-32}( \cos\frac{\pi +2k\pi}{5}+ i\sin \frac{\pi +2k\pi}{5})\), \(k=0,1,2,3,4\)}{\( x_k = \sqrt[5]{|-32|}( \cos \frac{\pi +k\pi}{5}+ i\sin \frac{\pi +k\pi}{5})\), \(k=0,1,2,3,4\)}{\( x_k = \sqrt[5]{|-32|}( \cos \frac{\pi +2k\pi}{5}+ \sin \frac{\pi +2k\pi}{5})\), \(k=0,1,2,3,4\)}{}{}}
\LANGcs{
\Question{
Vyberte správný vzorec pro řešení rovnice \(x^5 +32=0\).}
{\( x_k = \sqrt[5]{|-32|}( \cos\frac{\pi +2k\pi}{5}+ i\sin \frac{\pi +2k\pi}{5})\), \(k=0,1,2,3,4\)}{\( x_k = \sqrt[5]{-32}( \cos\frac{\pi +2k\pi}{5}+ i\sin \frac{\pi +2k\pi}{5})\), \(k=0,1,2,3,4\)}{\( x_k = \sqrt[5]{|-32|}( \cos \frac{\pi +k\pi}{5}+ i\sin \frac{\pi +k\pi}{5})\), \(k=0,1,2,3,4\)}{\( x_k = \sqrt[5]{|-32|}( \cos \frac{\pi +2k\pi}{5}+ \sin \frac{\pi +2k\pi}{5})\), \(k=0,1,2,3,4\)}{}{}}
\LANGsk{
\Question{
Vyberte správný vzorec pre riešenie rovnice \(x^5 +32=0\).}
{\( x_k = \sqrt[5]{|-32|}( \cos\frac{\pi +2k\pi}{5}+ i\sin \frac{\pi +2k\pi}{5})\), \(k=0,1,2,3,4\)}{\( x_k = \sqrt[5]{-32}( \cos\frac{\pi +2k\pi}{5}+ i\sin \frac{\pi +2k\pi}{5})\), \(k=0,1,2,3,4\)}{\( x_k = \sqrt[5]{|-32|}( \cos \frac{\pi +k\pi}{5}+ i\sin \frac{\pi +k\pi}{5})\), \(k=0,1,2,3,4\)}{\( x_k = \sqrt[5]{|-32|}( \cos \frac{\pi +2k\pi}{5}+ \sin \frac{\pi +2k\pi}{5})\), \(k=0,1,2,3,4\)}{}{}}
\LANGes{
\Question{
Determina la fórmula correcta para resolver la ecuación \(x^5 +32=0\)}
{\( x_k = \sqrt[5]{|-32|}( \cos\frac{\pi +2k\pi}{5}+ i\sin \frac{\pi +2k\pi}{5})\), \(k=0,1,2,3,4\)}{\( x_k = \sqrt[5]{-32}( \cos\frac{\pi +2k\pi}{5}+ i\sin \frac{\pi +2k\pi}{5})\), \(k=0,1,2,3,4\)}{\( x_k = \sqrt[5]{|-32|}( \cos \frac{\pi +k\pi}{5}+ i\sin \frac{\pi +k\pi}{5})\), \(k=0,1,2,3,4\)}{\( x_k = \sqrt[5]{|-32|}( \cos \frac{\pi +2k\pi}{5}+ \sin \frac{\pi +2k\pi}{5})\), \(k=0,1,2,3,4\)}{}{}}
\End

\Data\ID{2000002609}
\Updated{2021-12-01 04:12:35}
\Level{A}
\SubArea{Binomial equations}
\LANGen{
\Question{
The equation \(x^3 +27=0\) can be solved by factoring. The expression on the left can be factored as:}
{\( (x+3)(x^2-3x+9) \)}{\( (x+3)(x^2+3x+9) \)}{\( (x+3)^3 \)}{\( (x-3)(x^2+3x+9) \)}{}{}}
\LANGpl{
\Question{
Równanie \(x^3 +27=0\)     można rozwiązać za pomocą rozkładu na czynniki. Rozkład na czynniki ma postać:}
{\( (x+3)(x^2-3x+9) \)}{\( (x+3)(x^2+3x+9) \)}{\( (x+3)^3 \)}{\( (x-3)(x^2+3x+9) \)}{}{}}
\LANGcs{
\Question{
Rovnici \(x^3 +27=0\) lze řešit rozkladem na součin. Výraz na levé straně rovnice lze rozložit na:}
{\( (x+3)(x^2-3x+9) \)}{\( (x+3)(x^2+3x+9) \)}{\( (x+3)^3 \)}{\( (x-3)(x^2+3x+9) \)}{}{}}
\LANGsk{
\Question{
Rovnicu \(x^3 +27=0\) môžeme riešiť rozkladom na súčin. Výraz na ľavej strane rovnice rozložíme na:}
{\( (x+3)(x^2-3x+9) \)}{\( (x+3)(x^2+3x+9) \)}{\( (x+3)^3 \)}{\( (x-3)(x^2+3x+9) \)}{}{}}
\LANGes{
\Question{
La ecuación \(x^3 +27=0\) se puede resolver factorizando. La expresión de la izquierda se puede factorizar como:}
{\( (x+3)(x^2-3x+9) \)}{\( (x+3)(x^2+3x+9) \)}{\( (x+3)^3 \)}{\( (x-3)(x^2+3x+9) \)}{}{}}
\End

\Data\ID{2000002610}
\Updated{2021-12-01 04:12:40}
\Level{A}
\SubArea{Binomial equations}
\LANGen{
\Question{
One solution of the equation \( x^2-72i=0\) is \(x_1 = 6(1+i)\). Find the missing solution.}
{\( x_2 = -6(1+i) \)}{\( x_2 = 6(1-i) \)}{\( x_2 = 6(-1+i) \)}{\( x_2 = -6(1-i) \)}{}{}}
\LANGpl{
\Question{
Jednym z rozwiązań równania  \( x^2-72i=0\) jest \(x_1 = 6(1+i)\). Znajdź brakujące rozwiązanie.}
{\( x_2 = -6(1+i) \)}{\( x_2 = 6(1-i) \)}{\( x_2 = 6(-1+i) \)}{\( x_2 = -6(1-i) \)}{}{}}
\LANGcs{
\Question{
Jedním z řešení \( x^2-72i=0\) je \(x_1 = 6(1+i)\). Určete chybějící řešení.}
{\( x_2 = -6(1+i) \)}{\( x_2 = 6(1-i) \)}{\( x_2 = 6(-1+i) \)}{\( x_2 = -6(1-i) \)}{}{}}
\LANGsk{
\Question{
Jeden z koreňov rovnice \( x^2-72i=0\) je \(x_1 = 6(1+i)\). Určte chýbajúce riešenie.}
{\( x_2 = -6(1+i) \)}{\( x_2 = 6(1-i) \)}{\( x_2 = 6(-1+i) \)}{\( x_2 = -6(1-i) \)}{}{}}
\LANGes{
\Question{
Una solución de la ecuación \( x^2-72i=0\) es \(x_1 = 6(1+i)\). Determina la solución que falta.}
{\( x_2 = -6(1+i) \)}{\( x_2 = 6(1-i) \)}{\( x_2 = 6(-1+i) \)}{\( x_2 = -6(1-i) \)}{}{}}
\End

\Data\ID{2010013401}
\Updated{2022-11-02 09:11:59}
\Level{A}
\SubArea{Binomial equations}
\LANGen{
\Question{
Consider the equation \( x^4 + 16 = 0 \). Find the sum of all its roots in the set of complex numbers.}
{\( 0 \)}{\( 2 \)}{\( -2 \)}{\( 16 \)}{\( -16 \)}{}}
\LANGpl{
\Question{
Rozwiąż równanie w zbiorze liczb zespolonych \( x^4 + 16 = 0 \). Znajdź sumę wszystkich jego pierwiastków.}
{\( 0 \)}{\( 2 \)}{\( -2 \)}{\( 16 \)}{\( -16 \)}{}}
\LANGcs{
\Question{
Určete součet komplexních kořenů rovnice \( x^4 + 16 = 0 \).}
{\( 0 \)}{\( 2 \)}{\( -2 \)}{\( 16 \)}{\( -16 \)}{}}
\LANGsk{
\Question{
Súčet koreňov binomickej rovnice \( x^4 + 16 = 0 \) je:}
{\( 0 \)}{\( 2 \)}{\( -2 \)}{\( 16 \)}{\( -16 \)}{}}
\LANGes{
\Question{
Se considera la ecuación \( x^4 + 16 = 0 \). Halla la suma de todas sus raíces en el conjunto de los números complejos.}
{\( 0 \)}{\( 2 \)}{\( -2 \)}{\( 16 \)}{\( -16 \)}{}}
\End

\Data\ID{2010013402}
\Updated{2022-11-02 09:11:59}
\Level{A}
\SubArea{Binomial equations}
\LANGen{
\Question{
Solve the following equation in the set of complex numbers. (Solve the equation by substitution.)
\[ (3x + 2)^4 - 81 = 0 \]}
{\( \left\{-\frac53;\frac13;-\frac23+\mathrm{i} ;-\frac23-\mathrm{i} \right\} \)}{\( \left\{-\frac53;\frac13;\frac23+\mathrm{i} ;\frac23-\mathrm{i} \right\} \)}{\( \left\{\frac53;-\frac13;-\frac23+\mathrm{i} ;-\frac23-\mathrm{i} \right\} \)}{\( \left\{\frac53;-\frac13;\frac23+\mathrm{i} ;\frac23-\mathrm{i} \right\} \)}{}{}}
\LANGpl{
\Question{
Rozwiąż następujące równanie w zbiorze liczb zespolonych. (Rozwiąż równanie przez podstawienie.)
\[ (3x + 2)^4 - 81 = 0 \]}
{\( \left\{-\frac53;\frac13;-\frac23+\mathrm{i} ;-\frac23-\mathrm{i} \right\} \)}{\( \left\{-\frac53;\frac13;\frac23+\mathrm{i} ;\frac23-\mathrm{i} \right\} \)}{\( \left\{\frac53;-\frac13;-\frac23+\mathrm{i} ;-\frac23-\mathrm{i} \right\} \)}{\( \left\{\frac53;-\frac13;\frac23+\mathrm{i} ;\frac23-\mathrm{i} \right\} \)}{}{}}
\LANGcs{
\Question{
Určete množinu komplexních kořenů dané rovnice. (Použijte substituci.)
\[ (3x + 2)^4 - 81 = 0 \]}
{\( \left\{-\frac53;\frac13;-\frac23+\mathrm{i} ;-\frac23-\mathrm{i} \right\} \)}{\( \left\{-\frac53;\frac13;\frac23+\mathrm{i} ;\frac23-\mathrm{i} \right\} \)}{\( \left\{\frac53;-\frac13;-\frac23+\mathrm{i} ;-\frac23-\mathrm{i} \right\} \)}{\( \left\{\frac53;-\frac13;\frac23+\mathrm{i} ;\frac23-\mathrm{i} \right\} \)}{}{}}
\LANGsk{
\Question{
Nájdite množinu komplexných koreňov binomickej rovnice 
\[ (3x + 2)^4 - 81 = 0. \]
(Riešte pomocou substitúcie.)}
{\( \left\{-\frac53;\frac13;-\frac23+\mathrm{i} ;-\frac23-\mathrm{i} \right\} \)}{\( \left\{-\frac53;\frac13;\frac23+\mathrm{i} ;\frac23-\mathrm{i} \right\} \)}{\( \left\{\frac53;-\frac13;-\frac23+\mathrm{i} ;-\frac23-\mathrm{i} \right\} \)}{\( \left\{\frac53;-\frac13;\frac23+\mathrm{i} ;\frac23-\mathrm{i} \right\} \)}{}{}}
\LANGes{
\Question{
Resuelve la siguiente ecuación en el conjunto de los números complejos. (Utiliza un cambio de variable.)
\[ (3x + 2)^4 - 81 = 0 \]}
{\( \left\{-\frac53;\frac13;-\frac23+\mathrm{i} ;-\frac23-\mathrm{i} \right\} \)}{\( \left\{-\frac53;\frac13;\frac23+\mathrm{i} ;\frac23-\mathrm{i} \right\} \)}{\( \left\{\frac53;-\frac13;-\frac23+\mathrm{i} ;-\frac23-\mathrm{i} \right\} \)}{\( \left\{\frac53;-\frac13;\frac23+\mathrm{i} ;\frac23-\mathrm{i} \right\} \)}{}{}}
\End

\Data\ID{2010013403}
\Updated{2022-11-02 09:11:59}
\Level{A}
\SubArea{Binomial equations}
\LANGen{
\Question{
All the solutions of the equation \( x^6+3\sqrt5-6\mathrm{i} = 0 \) can be  displayed as points in the  rectangular coordinate system. What is the distance of the two most distant points?}
{\( 2\sqrt[3]3 \)}{\( 2\sqrt3 \)}{\( \sqrt3 \)}{\( \sqrt[3]9\)}{\( 2\sqrt[3]9\)}{}}
\LANGpl{
\Question{
Wszystkie rozwiązania równania \( x^6+3\sqrt5-6\mathrm{i} = 0 \) mogą być przedstawione, jako punkty w prostokątnym układzie współrzędnych. Jaka jest odległość dwóch najbardziej odległych punktów?}
{\( 2\sqrt[3]3 \)}{\( 2\sqrt3 \)}{\( \sqrt3 \)}{\( \sqrt[3]9\)}{\( 2\sqrt[3]9\)}{}}
\LANGcs{
\Question{
Všechna řešení rovnice \( x^6+3\sqrt5-6\mathrm{i} = 0 \) mohou být zobrazena v pravúhlém souřadnicovém systému. Jaká je vzdálenost nejvzdálenějších bodů?}
{\( 2\sqrt[3]3 \)}{\( 2\sqrt3 \)}{\( \sqrt3 \)}{\( \sqrt[3]9\)}{\( 2\sqrt[3]9\)}{}}
\LANGsk{
\Question{
Vzdialenosť dvoch najvzdialenejších obrazov koreňov binomickej rovnice \( x^6+3\sqrt5-6\mathrm{i} = 0 \) zobrazených v Gaussovej rovine, je:}
{\( 2\sqrt[3]3 \)}{\( 2\sqrt3 \)}{\( \sqrt3 \)}{\( \sqrt[3]9\)}{\( 2\sqrt[3]9\)}{}}
\LANGes{
\Question{
Todas las soluciones de la ecuación \( x^6+3\sqrt5-6\mathrm{i} = 0 \) se pueden representar en el plano cartesiano. ¿Cuál es la distancia entre los dos puntos más lejanos?}
{\( 2\sqrt[3]3 \)}{\( 2\sqrt3 \)}{\( \sqrt3 \)}{\( \sqrt[3]9\)}{\( 2\sqrt[3]9\)}{}}
\End

\Data\ID{2010013404}
\Updated{2022-11-02 09:11:59}
\Level{A}
\SubArea{Binomial equations}
\LANGen{
\Question{
Find the solution set of the following equation in the set of complex numbers.
\[ 
 x^{3} - 8 = 0
\]}
{\(\left\{2;\ -1 - \mathrm{i}\sqrt{3};\ -1 +\mathrm{i}\sqrt{3}\right\}\)}{\(\left\{2;\ 1 - \mathrm{i}\sqrt{3};\ 1 +\mathrm{i}\sqrt{3}\right\}\)}{\(\left\{2;\ 1 - \mathrm{i}\sqrt{3}\right\}\)}{\(\left\{2;\ - 1 + \mathrm{i}\sqrt{3}\right\}\)}{}{}}
\LANGpl{
\Question{
Znajdź zbiór rozwiązań następującego równania w zbiorze liczb zespolonych.
\[ 
 x^{3} - 8 = 0
\]}
{\(\left\{2;\ -1 - \mathrm{i}\sqrt{3};\ -1 +\mathrm{i}\sqrt{3}\right\}\)}{\(\left\{2;\ 1 - \mathrm{i}\sqrt{3};\ 1 +\mathrm{i}\sqrt{3}\right\}\)}{\(\left\{2;\ 1 - \mathrm{i}\sqrt{3}\right\}\)}{\(\left\{2;\ - 1 + \mathrm{i}\sqrt{3}\right\}\)}{}{}}
\LANGcs{
\Question{
Určete množinu všech komplexních kořenů dané rovnice.
\[ 
 x^{3} - 8 = 0
\]}
{\(\left\{2;\ -1 - \mathrm{i}\sqrt{3};\ -1 +\mathrm{i}\sqrt{3}\right\}\)}{\(\left\{2;\ 1 - \mathrm{i}\sqrt{3};\ 1 +\mathrm{i}\sqrt{3}\right\}\)}{\(\left\{2;\ 1 - \mathrm{i}\sqrt{3}\right\}\)}{\(\left\{2;\ - 1 + \mathrm{i}\sqrt{3}\right\}\)}{}{}}
\LANGsk{
\Question{
Nájdite množinu všetkých riešení danej rovnice v množine komplexných čísel.
\[ 
 x^{3} - 8 = 0
\]}
{\(\left\{2;\ -1 - \mathrm{i}\sqrt{3};\ -1 +\mathrm{i}\sqrt{3}\right\}\)}{\(\left\{2;\ 1 - \mathrm{i}\sqrt{3};\ 1 +\mathrm{i}\sqrt{3}\right\}\)}{\(\left\{2;\ 1 - \mathrm{i}\sqrt{3}\right\}\)}{\(\left\{2;\ - 1 + \mathrm{i}\sqrt{3}\right\}\)}{}{}}
\LANGes{
\Question{
Halla todas las soluciones de la siguiente ecuación en el conjunto de los números complejos.
\[ 
 x^{3} - 8 = 0
\]}
{\(\left\{2;\ -1 - \mathrm{i}\sqrt{3};\ -1 +\mathrm{i}\sqrt{3}\right\}\)}{\(\left\{2;\ 1 - \mathrm{i}\sqrt{3};\ 1 +\mathrm{i}\sqrt{3}\right\}\)}{\(\left\{2;\ 1 - \mathrm{i}\sqrt{3}\right\}\)}{\(\left\{2;\ - 1 + \mathrm{i}\sqrt{3}\right\}\)}{}{}}
\End

\Data\ID{2010013405}
\Updated{2022-11-02 09:11:22}
\Level{A}
\SubArea{Binomial equations}
\LANGen{
\Question{
Find the solution set of the following equation in the set of complex numbers.
\[ 
 x^{3} + 27 = 0
\]}
{\(\left\{-3;\ \frac32 - \mathrm{i}\frac{3\sqrt3} {2} ;\ \frac32 +\mathrm{i}\frac{3\sqrt{3}} {2} \right\}\)}{\(\left\{-3;\ -\frac32 + \mathrm{i}\frac{3\sqrt3} {2} ;\ -\frac32 -\mathrm{i}\frac{3\sqrt3} {2} \right\}\)}{\(\left\{-3;\ \frac32 - \mathrm{i}\frac{\sqrt3} {2} ;\ \frac32 +\mathrm{i}\frac{\sqrt3} {2} \right\}\)}{\(\left\{-3;\ -\frac32 + \mathrm{i}\frac{\sqrt3} {2} ;\ -\frac32 -\mathrm{i}\frac{\sqrt3} {2} \right\}\)}{}{}}
\LANGpl{
\Question{
Znajdź zbiór rozwiązań nasępującego równania w zbiorze liczb zespolonych.
\[ 
 x^{3} + 27 = 0
\]}
{\(\left\{-3;\ \frac32 - \mathrm{i}\frac{3\sqrt3} {2} ;\ \frac32 +\mathrm{i}\frac{3\sqrt{3}} {2} \right\}\)}{\(\left\{-3;\ -\frac32 + \mathrm{i}\frac{3\sqrt3} {2} ;\ -\frac32 -\mathrm{i}\frac{3\sqrt3} {2} \right\}\)}{\(\left\{-3;\ \frac32 - \mathrm{i}\frac{\sqrt3} {2} ;\ \frac32 +\mathrm{i}\frac{\sqrt3} {2} \right\}\)}{\(\left\{-3;\ -\frac32 + \mathrm{i}\frac{\sqrt3} {2} ;\ -\frac32 -\mathrm{i}\frac{\sqrt3} {2} \right\}\)}{}{}}
\LANGcs{
\Question{
Určete množinu všech komplexních kořenů dané rovnice.
\[ 
 x^{3} + 27 = 0
\]}
{\(\left\{-3;\ \frac32 - \mathrm{i}\frac{3\sqrt3} {2} ;\ \frac32 +\mathrm{i}\frac{3\sqrt{3}} {2} \right\}\)}{\(\left\{-3;\ -\frac32 + \mathrm{i}\frac{3\sqrt3} {2} ;\ -\frac32 -\mathrm{i}\frac{3\sqrt3} {2} \right\}\)}{\(\left\{-3;\ \frac32 - \mathrm{i}\frac{\sqrt3} {2} ;\ \frac32 +\mathrm{i}\frac{\sqrt3} {2} \right\}\)}{\(\left\{-3;\ -\frac32 + \mathrm{i}\frac{\sqrt3} {2} ;\ -\frac32 -\mathrm{i}\frac{\sqrt3} {2} \right\}\)}{}{}}
\LANGsk{
\Question{
Nájdite množinu všetkých riešení danej rovnice v množine komplexných čísel.
\[ 
 x^{3} + 27 = 0
\]}
{\(\left\{-3;\ \frac32 - \mathrm{i}\frac{3\sqrt3} {2} ;\ \frac32 +\mathrm{i}\frac{3\sqrt{3}} {2} \right\}\)}{\(\left\{-3;\ -\frac32 + \mathrm{i}\frac{3\sqrt3} {2} ;\ -\frac32 -\mathrm{i}\frac{3\sqrt3} {2} \right\}\)}{\(\left\{-3;\ \frac32 - \mathrm{i}\frac{\sqrt3} {2} ;\ \frac32 +\mathrm{i}\frac{\sqrt3} {2} \right\}\)}{\(\left\{-3;\ -\frac32 + \mathrm{i}\frac{\sqrt3} {2} ;\ -\frac32 -\mathrm{i}\frac{\sqrt3} {2} \right\}\)}{}{}}
\LANGes{
\Question{
Calcula todas las soluciones de la siguiente ecuación en el conjunto de los números complejos.
\[ 
 x^{3} + 27 = 0
\]}
{\(\left\{-3;\ \frac32 - \mathrm{i}\frac{3\sqrt3} {2} ;\ \frac32 +\mathrm{i}\frac{3\sqrt{3}} {2} \right\}\)}{\(\left\{-3;\ -\frac32 + \mathrm{i}\frac{3\sqrt3} {2} ;\ -\frac32 -\mathrm{i}\frac{3\sqrt3} {2} \right\}\)}{\(\left\{-3;\ \frac32 - \mathrm{i}\frac{\sqrt3} {2} ;\ \frac32 +\mathrm{i}\frac{\sqrt3} {2} \right\}\)}{\(\left\{-3;\ -\frac32 + \mathrm{i}\frac{\sqrt3} {2} ;\ -\frac32 -\mathrm{i}\frac{\sqrt3} {2} \right\}\)}{}{}}
\End

\Data\ID{2010013406}
\Updated{2022-11-02 09:11:22}
\Level{A}
\SubArea{Binomial equations}
\LANGen{
\Question{
Find the solution set of the following equation in the set of complex numbers.
\[ 
 x^{3} + 8\mathrm{i} = 0
\]}
{\(\left\{2\mathrm{i};\ \sqrt{3} -\mathrm{i};\ -\sqrt{3}-\mathrm{i}\right\}\)}{\(\left\{ -2\mathrm{i};\ \sqrt{3} -\mathrm{i};\ -\sqrt{3}-\mathrm{i}\right\}\)}{\(\left\{ -2;\ -\sqrt{3} +\mathrm{i};\ \sqrt{3}+\mathrm{i}\right\}\)}{\(\left\{ 2;\ -\sqrt{3} +\mathrm{i};\ \sqrt{3}+\mathrm{i}\right\}\)}{}{}}
\LANGpl{
\Question{
Znajdź zbiór rozwiązań następującego równania w zbiorze liczb zespolonych. 
\[ 
 x^{3} + 8\mathrm{i} = 0
\]}
{\(\left\{2\mathrm{i};\ \sqrt{3} -\mathrm{i};\ -\sqrt{3}-\mathrm{i}\right\}\)}{\(\left\{ -2\mathrm{i};\ \sqrt{3} -\mathrm{i};\ -\sqrt{3}-\mathrm{i}\right\}\)}{\(\left\{ -2;\ -\sqrt{3} +\mathrm{i};\ \sqrt{3}+\mathrm{i}\right\}\)}{\(\left\{ 2;\ -\sqrt{3} +\mathrm{i};\ \sqrt{3}+\mathrm{i}\right\}\)}{}{}}
\LANGcs{
\Question{
Určete množinu všech komplexních kořenů dané kvadratické rovnice.
\[ 
 x^{3} + 8\mathrm{i} = 0
\]}
{\(\left\{2\mathrm{i};\ \sqrt{3} -\mathrm{i};\ -\sqrt{3}-\mathrm{i}\right\}\)}{\(\left\{ -2\mathrm{i};\ \sqrt{3} -\mathrm{i};\ -\sqrt{3}-\mathrm{i}\right\}\)}{\(\left\{ -2;\ -\sqrt{3} +\mathrm{i};\ \sqrt{3}+\mathrm{i}\right\}\)}{\(\left\{ 2;\ -\sqrt{3} +\mathrm{i};\ \sqrt{3}+\mathrm{i}\right\}\)}{}{}}
\LANGsk{
\Question{
Nájdite množinu všetkých riešení danej rovnice v množine komplexných čísel.
\[ 
 x^{3} + 8\mathrm{i} = 0
\]}
{\(\left\{2\mathrm{i};\ \sqrt{3} -\mathrm{i};\ -\sqrt{3}-\mathrm{i}\right\}\)}{\(\left\{ -2\mathrm{i};\ \sqrt{3} -\mathrm{i};\ -\sqrt{3}-\mathrm{i}\right\}\)}{\(\left\{ -2;\ -\sqrt{3} +\mathrm{i};\ \sqrt{3}+\mathrm{i}\right\}\)}{\(\left\{ 2;\ -\sqrt{3} +\mathrm{i};\ \sqrt{3}+\mathrm{i}\right\}\)}{}{}}
\LANGes{
\Question{
Halla todas las soluciones de la siguiente ecuación en el conjunto de los números complejos.
\[ 
 x^{3} + 8\mathrm{i} = 0
\]}
{\(\left\{2\mathrm{i};\ \sqrt{3} -\mathrm{i};\ -\sqrt{3}-\mathrm{i}\right\}\)}{\(\left\{ -2\mathrm{i};\ \sqrt{3} -\mathrm{i};\ -\sqrt{3}-\mathrm{i}\right\}\)}{\(\left\{ -2;\ -\sqrt{3} +\mathrm{i};\ \sqrt{3}+\mathrm{i}\right\}\)}{\(\left\{ 2;\ -\sqrt{3} +\mathrm{i};\ \sqrt{3}+\mathrm{i}\right\}\)}{}{}}
\End

\Data\ID{2010013407}
\Updated{2022-11-02 09:11:22}
\Level{A}
\SubArea{Binomial equations}
\LANGen{
\Question{
Two solutions of the equation 
\[ 
 x^{3} + 1 - \mathrm{i} = 0
\]
are 
\[ 
\begin{aligned}x_{1}& = \root{6}\of{2}\left (\cos \frac{\pi}
{4} + \mathrm{i}\sin \frac{\pi}
{4} \right ),&
\\x_{2}& = \root{6}\of{2}\left (\cos \frac{11}
{12}\pi + \mathrm{i}\sin \frac{11}
{12}\pi \right ).
\\ \end{aligned}
\]
Find the third solution.}
{\(x_{3} = \root{6}\of{2}\left (\cos \frac{19}
{12}\pi + \mathrm{i}\sin \frac{19}
{12}\pi \right )\)}{\(x_{3} = \root{6}\of{2}\left (\cos \frac{7}
{12}\pi + \mathrm{i}\sin \frac{7}
{12}\pi \right )\)}{\(x_{3} = \root{6}\of{2}\left (\cos \frac{5}
{12}\pi + \mathrm{i}\sin \frac{5}
{12}\pi \right )\)}{\(x_{3} = \root{6}\of{2}\left (\cos \frac{13}
{12}\pi + \mathrm{i}\sin \frac{13}
{12}\pi \right )\)}{}{}}
\LANGpl{
\Question{
Dwa rozwiązania równania  
\[ 
 x^{3} + 1 - \mathrm{i} = 0
\]
są równe:
\[ 
\begin{aligned}x_{1}& = \root{6}\of{2}\left (\cos \frac{\pi}
{4} + \mathrm{i}\sin \frac{\pi}
{4} \right ),&
\\x_{2}& = \root{6}\of{2}\left (\cos \frac{11}
{12}\pi + \mathrm{i}\sin \frac{11}
{12}\pi \right ).
\\ \end{aligned}
\]
Znajdź trzecie rozwiązanie.}
{\(x_{3} = \root{6}\of{2}\left (\cos \frac{19}
{12}\pi + \mathrm{i}\sin \frac{19}
{12}\pi \right )\)}{\(x_{3} = \root{6}\of{2}\left (\cos \frac{7}
{12}\pi + \mathrm{i}\sin \frac{7}
{12}\pi \right )\)}{\(x_{3} = \root{6}\of{2}\left (\cos \frac{5}
{12}\pi + \mathrm{i}\sin \frac{5}
{12}\pi \right )\)}{\(x_{3} = \root{6}\of{2}\left (\cos \frac{13}
{12}\pi + \mathrm{i}\sin \frac{13}
{12}\pi \right )\)}{}{}}
\LANGcs{
\Question{
Dva kořeny rovnice
\[ 
 x^{3} + 1 - \mathrm{i} = 0
\]
jsou
\[ 
\begin{aligned}x_{1}& = \root{6}\of{2}\left (\cos \frac{\pi}
{4} + \mathrm{i}\sin \frac{\pi}
{4} \right ),&
\\x_{2}& = \root{6}\of{2}\left (\cos \frac{11}
{12}\pi + \mathrm{i}\sin \frac{11}
{12}\pi \right ).
\\ \end{aligned}
\]
Určete třetí kořen.}
{\(x_{3} = \root{6}\of{2}\left (\cos \frac{19}
{12}\pi + \mathrm{i}\sin \frac{19}
{12}\pi \right )\)}{\(x_{3} = \root{6}\of{2}\left (\cos \frac{7}
{12}\pi + \mathrm{i}\sin \frac{7}
{12}\pi \right )\)}{\(x_{3} = \root{6}\of{2}\left (\cos \frac{5}
{12}\pi + \mathrm{i}\sin \frac{5}
{12}\pi \right )\)}{\(x_{3} = \root{6}\of{2}\left (\cos \frac{13}
{12}\pi + \mathrm{i}\sin \frac{13}
{12}\pi \right )\)}{}{}}
\LANGsk{
\Question{
Dve riešenia rovnice
\[ 
 x^{3} + 1 - \mathrm{i} = 0
\]
sú 
\[ 
\begin{aligned}x_{1}& = \root{6}\of{2}\left (\cos \frac{\pi}
{4} + \mathrm{i}\sin \frac{\pi}
{4} \right ),&
\\x_{2}& = \root{6}\of{2}\left (\cos \frac{11}
{12}\pi + \mathrm{i}\sin \frac{11}
{12}\pi \right ).
\\ \end{aligned}
\]
Nájdite tretie riešenie.}
{\(x_{3} = \root{6}\of{2}\left (\cos \frac{19}
{12}\pi + \mathrm{i}\sin \frac{19}
{12}\pi \right )\)}{\(x_{3} = \root{6}\of{2}\left (\cos \frac{7}
{12}\pi + \mathrm{i}\sin \frac{7}
{12}\pi \right )\)}{\(x_{3} = \root{6}\of{2}\left (\cos \frac{5}
{12}\pi + \mathrm{i}\sin \frac{5}
{12}\pi \right )\)}{\(x_{3} = \root{6}\of{2}\left (\cos \frac{13}
{12}\pi + \mathrm{i}\sin \frac{13}
{12}\pi \right )\)}{}{}}
\LANGes{
\Question{
Dos soluciones de la ecuación 
\[ 
 x^{3} + 1 - \mathrm{i} = 0
\]
son 
\[ 
\begin{aligned}x_{1}& = \root{6}\of{2}\left (\cos \frac{\pi}
{4} + \mathrm{i}\sin \frac{\pi}
{4} \right ),&
\\x_{2}& = \root{6}\of{2}\left (\cos \frac{11}
{12}\pi + \mathrm{i}\sin \frac{11}
{12}\pi \right ).
\\ \end{aligned}
\]
Calcula la tercera solución.}
{\(x_{3} = \root{6}\of{2}\left (\cos \frac{19}
{12}\pi + \mathrm{i}\sin \frac{19}
{12}\pi \right )\)}{\(x_{3} = \root{6}\of{2}\left (\cos \frac{7}
{12}\pi + \mathrm{i}\sin \frac{7}
{12}\pi \right )\)}{\(x_{3} = \root{6}\of{2}\left (\cos \frac{5}
{12}\pi + \mathrm{i}\sin \frac{5}
{12}\pi \right )\)}{\(x_{3} = \root{6}\of{2}\left (\cos \frac{13}
{12}\pi + \mathrm{i}\sin \frac{13}
{12}\pi \right )\)}{}{}}
\End

\Data\ID{2010013408}
\Updated{2022-11-02 09:11:22}
\Level{A}
\SubArea{Binomial equations}
\LANGen{
\Question{
Three solutions of the equation 
\[ 
 x^{4} - 2\mathrm{i} = 0
\]
are 
 \[x_{1} = \root{4}\of{2}\left (\cos \frac{1}{8}\pi + \mathrm{i}\sin \frac{1}{8}\pi   \right ), \ x_{2} = \root{4}\of{2}\left (\cos \frac{5}{8}\pi + \mathrm{i}\sin \frac{5}{8}\pi   \right ),\  x_{3} = \root{4}\of{2}\left (\cos \frac{9}{8}\pi + \mathrm{i}\sin \frac{9}{8}\pi   \right ).\]
Find the fourth solution.}
{\(x_{4} = \root{4}\of{2}\left (\cos \frac{13}{8}\pi + \mathrm{i}\sin \frac{13}{8}\pi   \right )\)}{\(x_{4} = \root{4}\of{2}\left (\cos \frac{11}{8}\pi + \mathrm{i}\sin \frac{11}{8}\pi   \right )\)}{\(x_{4} = \root{4}\of{2}\left (\cos \frac{15}{8}\pi + \mathrm{i}\sin \frac{15}{8}\pi   \right )\)}{\(x_{4} = \root{4}\of{2}\left (\cos \frac{3}{8}\pi + \mathrm{i}\sin \frac{3}{8}\pi   \right )\)}{}{}}
\LANGpl{
\Question{
Trzy rozwiązania równania 
\[ 
 x^{4} - 2\mathrm{i} = 0
\]
są równe:
 \[x_{1} = \root{4}\of{2}\left (\cos \frac{1}{8}\pi + \mathrm{i}\sin \frac{1}{8}\pi   \right ), \ x_{2} = \root{4}\of{2}\left (\cos \frac{5}{8}\pi + \mathrm{i}\sin \frac{5}{8}\pi   \right ),\  x_{3} = \root{4}\of{2}\left (\cos \frac{9}{8}\pi + \mathrm{i}\sin \frac{9}{8}\pi   \right ).\]
Znajdź czwarte rozwiązanie.}
{\(x_{4} = \root{4}\of{2}\left (\cos \frac{13}{8}\pi + \mathrm{i}\sin \frac{13}{8}\pi   \right )\)}{\(x_{4} = \root{4}\of{2}\left (\cos \frac{11}{8}\pi + \mathrm{i}\sin \frac{11}{8}\pi   \right )\)}{\(x_{4} = \root{4}\of{2}\left (\cos \frac{15}{8}\pi + \mathrm{i}\sin \frac{15}{8}\pi   \right )\)}{\(x_{4} = \root{4}\of{2}\left (\cos \frac{3}{8}\pi + \mathrm{i}\sin \frac{3}{8}\pi   \right )\)}{}{}}
\LANGcs{
\Question{
Tři kořeny rovnice
\[ 
 x^{4} - 2\mathrm{i} = 0
\]
jsou
 \[x_{1} = \root{4}\of{2}\left (\cos \frac{1}{8}\pi + \mathrm{i}\sin \frac{1}{8}\pi   \right ), \ x_{2} = \root{4}\of{2}\left (\cos \frac{5}{8}\pi + \mathrm{i}\sin \frac{5}{8}\pi   \right ),\  x_{3} = \root{4}\of{2}\left (\cos \frac{9}{8}\pi + \mathrm{i}\sin \frac{9}{8}\pi   \right ).\]
Určete čtvrtý kořen.}
{\(x_{4} = \root{4}\of{2}\left (\cos \frac{13}{8}\pi + \mathrm{i}\sin \frac{13}{8}\pi   \right )\)}{\(x_{4} = \root{4}\of{2}\left (\cos \frac{11}{8}\pi + \mathrm{i}\sin \frac{11}{8}\pi   \right )\)}{\(x_{4} = \root{4}\of{2}\left (\cos \frac{15}{8}\pi + \mathrm{i}\sin \frac{15}{8}\pi   \right )\)}{\(x_{4} = \root{4}\of{2}\left (\cos \frac{3}{8}\pi + \mathrm{i}\sin \frac{3}{8}\pi   \right )\)}{}{}}
\LANGsk{
\Question{
Tri riešenia rovnice
\[ 
 x^{4} - 2\mathrm{i} = 0
\]
sú 
 \[x_{1} = \root{4}\of{2}\left (\cos \frac{1}{8}\pi + \mathrm{i}\sin \frac{1}{8}\pi   \right ), \ x_{2} = \root{4}\of{2}\left (\cos \frac{5}{8}\pi + \mathrm{i}\sin \frac{5}{8}\pi   \right ),\  x_{3} = \root{4}\of{2}\left (\cos \frac{9}{8}\pi + \mathrm{i}\sin \frac{9}{8}\pi   \right ).\]
Nájdite štvrté riešenie.}
{\(x_{4} = \root{4}\of{2}\left (\cos \frac{13}{8}\pi + \mathrm{i}\sin \frac{13}{8}\pi   \right )\)}{\(x_{4} = \root{4}\of{2}\left (\cos \frac{11}{8}\pi + \mathrm{i}\sin \frac{11}{8}\pi   \right )\)}{\(x_{4} = \root{4}\of{2}\left (\cos \frac{15}{8}\pi + \mathrm{i}\sin \frac{15}{8}\pi   \right )\)}{\(x_{4} = \root{4}\of{2}\left (\cos \frac{3}{8}\pi + \mathrm{i}\sin \frac{3}{8}\pi   \right )\)}{}{}}
\LANGes{
\Question{
Tres soluciones de la ecuación 
\[ 
 x^{4} - 2\mathrm{i} = 0
\]
son 
 \[x_{1} = \root{4}\of{2}\left (\cos \frac{1}{8}\pi + \mathrm{i}\sin \frac{1}{8}\pi   \right ), \ x_{2} = \root{4}\of{2}\left (\cos \frac{5}{8}\pi + \mathrm{i}\sin \frac{5}{8}\pi   \right ),\  x_{3} = \root{4}\of{2}\left (\cos \frac{9}{8}\pi + \mathrm{i}\sin \frac{9}{8}\pi   \right ).\]
Calcula la cuarta solución.}
{\(x_{4} = \root{4}\of{2}\left (\cos \frac{13}{8}\pi + \mathrm{i}\sin \frac{13}{8}\pi   \right )\)}{\(x_{4} = \root{4}\of{2}\left (\cos \frac{11}{8}\pi + \mathrm{i}\sin \frac{11}{8}\pi   \right )\)}{\(x_{4} = \root{4}\of{2}\left (\cos \frac{15}{8}\pi + \mathrm{i}\sin \frac{15}{8}\pi   \right )\)}{\(x_{4} = \root{4}\of{2}\left (\cos \frac{3}{8}\pi + \mathrm{i}\sin \frac{3}{8}\pi   \right )\)}{}{}}
\End

\Data\ID{2010013409}
\Updated{2022-11-02 09:11:22}
\Level{A}
\SubArea{Binomial equations}
\LANGen{
\Question{
Three solutions of the equation 
\[ 
 x^{4} + 8\mathrm{i} = 0
\]
are 
 \[x_{1} = \root{4}\of{8}\left (\cos \frac{3}{8}\pi + \mathrm{i}\sin \frac{3}{8}\pi   \right ), \ x_{2} = \root{4}\of{8}\left (\cos \frac{7}{8}\pi + \mathrm{i}\sin \frac{7}{8}\pi   \right ),\  x_{3} = \root{4}\of{8}\left (\cos \frac{15}{8}\pi + \mathrm{i}\sin \frac{15}{8}\pi   \right ).\]
Find the fourth solution.}
{\(x_{4} = \root{4}\of{8}\left (\cos \frac{11}{8}\pi + \mathrm{i}\sin \frac{11}{8}\pi   \right )\)}{\(x_{4} = \root{4}\of{8}\left (\cos \frac{9}{8}\pi + \mathrm{i}\sin \frac{9}{8}\pi   \right )\)}{\(x_{4} = \root{4}\of{8}\left (\cos \frac{5}{8}\pi + \mathrm{i}\sin \frac{5}{8}\pi   \right )\)}{\(x_{4} = \root{4}\of{8}\left (\cos \frac{1}{8}\pi + \mathrm{i}\sin \frac{1}{8}\pi   \right )\)}{}{}}
\LANGpl{
\Question{
Trzy rozwiązania równania 
\[ 
 x^{4} + 8\mathrm{i} = 0
\]
są równe:
 \[x_{1} = \root{4}\of{8}\left (\cos \frac{3}{8}\pi + \mathrm{i}\sin \frac{3}{8}\pi   \right ), \ x_{2} = \root{4}\of{8}\left (\cos \frac{7}{8}\pi + \mathrm{i}\sin \frac{7}{8}\pi   \right ),\  x_{3} = \root{4}\of{8}\left (\cos \frac{15}{8}\pi + \mathrm{i}\sin \frac{15}{8}\pi   \right ).\]
Znajdź czwarte rozwiązanie.}
{\(x_{4} = \root{4}\of{8}\left (\cos \frac{11}{8}\pi + \mathrm{i}\sin \frac{11}{8}\pi   \right )\)}{\(x_{4} = \root{4}\of{8}\left (\cos \frac{9}{8}\pi + \mathrm{i}\sin \frac{9}{8}\pi   \right )\)}{\(x_{4} = \root{4}\of{8}\left (\cos \frac{5}{8}\pi + \mathrm{i}\sin \frac{5}{8}\pi   \right )\)}{\(x_{4} = \root{4}\of{8}\left (\cos \frac{1}{8}\pi + \mathrm{i}\sin \frac{1}{8}\pi   \right )\)}{}{}}
\LANGcs{
\Question{
Tři kořeny rovnice
\[ 
 x^{4} + 8\mathrm{i} = 0
\]
jsou
 \[x_{1} = \root{4}\of{8}\left (\cos \frac{3}{8}\pi + \mathrm{i}\sin \frac{3}{8}\pi   \right ), \ x_{2} = \root{4}\of{8}\left (\cos \frac{7}{8}\pi + \mathrm{i}\sin \frac{7}{8}\pi   \right ),\  x_{3} = \root{4}\of{8}\left (\cos \frac{15}{8}\pi + \mathrm{i}\sin \frac{15}{8}\pi   \right ).\]
Určete čtvrtý kořen.}
{\(x_{4} = \root{4}\of{8}\left (\cos \frac{11}{8}\pi + \mathrm{i}\sin \frac{11}{8}\pi   \right )\)}{\(x_{4} = \root{4}\of{8}\left (\cos \frac{9}{8}\pi + \mathrm{i}\sin \frac{9}{8}\pi   \right )\)}{\(x_{4} = \root{4}\of{8}\left (\cos \frac{5}{8}\pi + \mathrm{i}\sin \frac{5}{8}\pi   \right )\)}{\(x_{4} = \root{4}\of{8}\left (\cos \frac{1}{8}\pi + \mathrm{i}\sin \frac{1}{8}\pi   \right )\)}{}{}}
\LANGsk{
\Question{
Tri riešenia rovnice
\[ 
 x^{4} + 8\mathrm{i} = 0
\]
sú 
 \[x_{1} = \root{4}\of{8}\left (\cos \frac{3}{8}\pi + \mathrm{i}\sin \frac{3}{8}\pi   \right ), \ x_{2} = \root{4}\of{8}\left (\cos \frac{7}{8}\pi + \mathrm{i}\sin \frac{7}{8}\pi   \right ),\  x_{3} = \root{4}\of{8}\left (\cos \frac{15}{8}\pi + \mathrm{i}\sin \frac{15}{8}\pi   \right ).\]
Nájdite štvrté riešenie.}
{\(x_{4} = \root{4}\of{8}\left (\cos \frac{11}{8}\pi + \mathrm{i}\sin \frac{11}{8}\pi   \right )\)}{\(x_{4} = \root{4}\of{8}\left (\cos \frac{9}{8}\pi + \mathrm{i}\sin \frac{9}{8}\pi   \right )\)}{\(x_{4} = \root{4}\of{8}\left (\cos \frac{5}{8}\pi + \mathrm{i}\sin \frac{5}{8}\pi   \right )\)}{\(x_{4} = \root{4}\of{8}\left (\cos \frac{1}{8}\pi + \mathrm{i}\sin \frac{1}{8}\pi   \right )\)}{}{}}
\LANGes{
\Question{
Tres soluciones de la ecuación 
\[ 
 x^{4} + 8\mathrm{i} = 0
\]
son 
 \[x_{1} = \root{4}\of{8}\left (\cos \frac{3}{8}\pi + \mathrm{i}\sin \frac{3}{8}\pi   \right ), \ x_{2} = \root{4}\of{8}\left (\cos \frac{7}{8}\pi + \mathrm{i}\sin \frac{7}{8}\pi   \right ),\  x_{3} = \root{4}\of{8}\left (\cos \frac{15}{8}\pi + \mathrm{i}\sin \frac{15}{8}\pi   \right ).\]
Calcula la cuarta solución.}
{\(x_{4} = \root{4}\of{8}\left (\cos \frac{11}{8}\pi + \mathrm{i}\sin \frac{11}{8}\pi   \right )\)}{\(x_{4} = \root{4}\of{8}\left (\cos \frac{9}{8}\pi + \mathrm{i}\sin \frac{9}{8}\pi   \right )\)}{\(x_{4} = \root{4}\of{8}\left (\cos \frac{5}{8}\pi + \mathrm{i}\sin \frac{5}{8}\pi   \right )\)}{\(x_{4} = \root{4}\of{8}\left (\cos \frac{1}{8}\pi + \mathrm{i}\sin \frac{1}{8}\pi   \right )\)}{}{}}
\End

\Data\ID{2010013410}
\Updated{2022-11-02 09:11:22}
\Level{A}
\SubArea{Binomial equations}
\LANGen{
\Question{
Consider the equation
\[ 
 x^{4} = ( 2 - \mathrm{i} )^{4}.
\]
Find the product of all its roots in the set of complex numbers.}
{\(7+24\mathrm{i}\)}{\(-7-24\mathrm{i}\)}{\(16\)}{\(-16\)}{}{}}
\LANGpl{
\Question{
Rozwiąż równanie w zbiorze liczb zespolonych
\[ 
 x^{4} = ( 2 - \mathrm{i} )^{4}.
\]
Znajdź iloczyn wszystkich jego pierwiastków.}
{\(7+24\mathrm{i}\)}{\(-7-24\mathrm{i}\)}{\(16\)}{\(-16\)}{}{}}
\LANGcs{
\Question{
Určete součin všech komplexních kořenů rovnice 
\[ 
 x^{4} = ( 2 - \mathrm{i} )^{4}.
\]}
{\(7+24\mathrm{i}\)}{\(-7-24\mathrm{i}\)}{\(16\)}{\(-16\)}{}{}}
\LANGsk{
\Question{
Určte súčin všetkých komplexných koreňov rovnice
\[ 
 x^{4} = ( 2 - \mathrm{i} )^{4}.
\]}
{\(7+24\mathrm{i}\)}{\(-7-24\mathrm{i}\)}{\(16\)}{\(-16\)}{}{}}
\LANGes{
\Question{
Dada al ecuación
\[ 
 x^{4} = ( 2 - \mathrm{i} )^{4}.
\]
Halla el producto de todas sus raíces en el conjunto de los números complejos.}
{\(7+24\mathrm{i}\)}{\(-7-24\mathrm{i}\)}{\(16\)}{\(-16\)}{}{}}
\End

\Data\ID{2010013411}
\Updated{2022-11-02 09:11:22}
\Level{A}
\SubArea{Binomial equations}
\LANGen{
\Question{
Consider the equation
\[ 
 x^{4} = ( 2 + \mathrm{i} )^{4}.
\]
Find the product of all its roots in the set of complex numbers.}
{\(7-24\mathrm{i}\)}{\(-7+24\mathrm{i}\)}{\(16\)}{\(-16\)}{}{}}
\LANGpl{
\Question{
Rozwiąż równanie w zbiorze liczb zespolonych
\[ 
 x^{4} = ( 2 + \mathrm{i} )^{4}.
\]
Znajdź iloczyn wszystkich jego pierwiastków.}
{\(7-24\mathrm{i}\)}{\(-7+24\mathrm{i}\)}{\(16\)}{\(-16\)}{}{}}
\LANGcs{
\Question{
Určete součin všech komplexních kořenů rovnice
\[ 
 x^{4} = ( 2 + \mathrm{i} )^{4}.
\]}
{\(7-24\mathrm{i}\)}{\(-7+24\mathrm{i}\)}{\(16\)}{\(-16\)}{}{}}
\LANGsk{
\Question{
Určte súčin všetkých komplexných koreňov rovnice
\[ 
 x^{4} = ( 2 + \mathrm{i} )^{4}.
\]}
{\(7-24\mathrm{i}\)}{\(-7+24\mathrm{i}\)}{\(16\)}{\(-16\)}{}{}}
\LANGes{
\Question{
Dada la ecuación
\[ 
 x^{4} = ( 2 + \mathrm{i} )^{4}.
\]
Halla el producto de todas sus raíces en el conjunto de los números complejos.}
{\(7-24\mathrm{i}\)}{\(-7+24\mathrm{i}\)}{\(16\)}{\(-16\)}{}{}}
\End

\Data\ID{2010013412}
\Updated{2022-11-02 09:11:22}
\Level{A}
\SubArea{Binomial equations}
\LANGen{
\Question{
Which of the given numbers does not belong to the solution set of the following equation?
\[x^{4}-1+\mathrm{i}=0\]}
{\(\root{8}\of{2}\left (\cos \frac{3\pi}{16} + \mathrm{i}\sin \frac{3\pi}{16}\right )\)}{\(-\root{4}\of{1-\mathrm{i}}\)}{\(-\mathrm{i}\root{4}\of{1-\mathrm{i}}\)}{\(\root{8}\of{2}\left (\cos \left(-\frac{\pi}{16}\right) + \mathrm{i}\sin \left(-\frac{\pi}{16}\right)\right )\)}{}{}}
\LANGpl{
\Question{
Która z liczb nie należy do zbioru rozwiązań poniższego równania?
\[x^{4}-1+\mathrm{i}=0\]}
{\(\root{8}\of{2}\left (\cos \frac{3\pi}{16} + \mathrm{i}\sin \frac{3\pi}{16}\right )\)}{\(-\root{4}\of{1-\mathrm{i}}\)}{\(-\mathrm{i}\root{4}\of{1-\mathrm{i}}\)}{\(\root{8}\of{2}\left (\cos \left(-\frac{\pi}{16}\right) + \mathrm{i}\sin \left(-\frac{\pi}{16}\right)\right )\)}{}{}}
\LANGcs{
\Question{
Které z čísel nepatří do množiny řešení dané rovnice?
\[x^{4}-1+\mathrm{i}=0\]}
{\(\root{8}\of{2}\left (\cos \frac{3\pi}{16} + \mathrm{i}\sin \frac{3\pi}{16}\right )\)}{\(-\root{4}\of{1-\mathrm{i}}\)}{\(-\mathrm{i}\root{4}\of{1-\mathrm{i}}\)}{\(\root{8}\of{2}\left (\cos \left(-\frac{\pi}{16}\right) + \mathrm{i}\sin \left(-\frac{\pi}{16}\right)\right )\)}{}{}}
\LANGsk{
\Question{
Ktoré z daných čísel nepatrí do množiny riešení danej rovnice?
\[x^{4}-1+\mathrm{i}=0\]}
{\(\root{8}\of{2}\left (\cos \frac{3\pi}{16} + \mathrm{i}\sin \frac{3\pi}{16}\right )\)}{\(-\root{4}\of{1-\mathrm{i}}\)}{\(-\mathrm{i}\root{4}\of{1-\mathrm{i}}\)}{\(\root{8}\of{2}\left (\cos \left(-\frac{\pi}{16}\right) + \mathrm{i}\sin \left(-\frac{\pi}{16}\right)\right )\)}{}{}}
\LANGes{
\Question{
¿Cuál de estos números no pertenece al conjunto de soluciones de la siguiente ecuación?
\[x^{4}-1+\mathrm{i}=0\]}
{\(\root{8}\of{2}\left (\cos \frac{3\pi}{16} + \mathrm{i}\sin \frac{3\pi}{16}\right )\)}{\(-\root{4}\of{1-\mathrm{i}}\)}{\(-\mathrm{i}\root{4}\of{1-\mathrm{i}}\)}{\(\root{8}\of{2}\left (\cos \left(-\frac{\pi}{16}\right) + \mathrm{i}\sin \left(-\frac{\pi}{16}\right)\right )\)}{}{}}
\End

\Data\ID{2010013413}
\Updated{2022-11-02 09:11:22}
\Level{A}
\SubArea{Binomial equations}
\LANGen{
\Question{
Which of the given numbers does not belong to the solution set of the following equation?
\[x^{4}+1+\mathrm{i}=0\]}
{\(\root{8}\of{2}\left (\cos \frac{3\pi}{16} + \mathrm{i}\sin \frac{3\pi}{16}\right )\)}{\(\mathrm{i}\root{4}\of{-1-\mathrm{i}}\)}{\(\root{4}\of{-1-\mathrm{i}}\)}{\(\root{8}\of{2}\left (\cos \frac{5\pi}{16} + \mathrm{i}\sin \frac{5\pi}{16}\right )\)}{}{}}
\LANGpl{
\Question{
Która z podanych liczb nie należy do zbioru rozwiązań poniższego równania?
\[x^{4}+1+\mathrm{i}=0\]}
{\(\root{8}\of{2}\left (\cos \frac{3\pi}{16} + \mathrm{i}\sin \frac{3\pi}{16}\right )\)}{\(\mathrm{i}\root{4}\of{-1-\mathrm{i}}\)}{\(\root{4}\of{-1-\mathrm{i}}\)}{\(\root{8}\of{2}\left (\cos \frac{5\pi}{16} + \mathrm{i}\sin \frac{5\pi}{16}\right )\)}{}{}}
\LANGcs{
\Question{
Které z čísel nepatří do množiny řešení dané rovnice?
\[x^{4}+1+\mathrm{i}=0\]}
{\(\root{8}\of{2}\left (\cos \frac{3\pi}{16} + \mathrm{i}\sin \frac{3\pi}{16}\right )\)}{\(\mathrm{i}\root{4}\of{-1-\mathrm{i}}\)}{\(\root{4}\of{-1-\mathrm{i}}\)}{\(\root{8}\of{2}\left (\cos \frac{5\pi}{16} + \mathrm{i}\sin \frac{5\pi}{16}\right )\)}{}{}}
\LANGsk{
\Question{
Ktoré z daných čísel nepatrí do množiny riešení danej rovnice?
\[x^{4}+1+\mathrm{i}=0\]}
{\(\root{8}\of{2}\left (\cos \frac{3\pi}{16} + \mathrm{i}\sin \frac{3\pi}{16}\right )\)}{\(\mathrm{i}\root{4}\of{-1-\mathrm{i}}\)}{\(\root{4}\of{-1-\mathrm{i}}\)}{\(\root{8}\of{2}\left (\cos \frac{5\pi}{16} + \mathrm{i}\sin \frac{5\pi}{16}\right )\)}{}{}}
\LANGes{
\Question{
¿Cuál de estos números no pertenece al conjunto de soluciones de la siguiente ecuación?
\[x^{4}+1+\mathrm{i}=0\]}
{\(\root{8}\of{2}\left (\cos \frac{3\pi}{16} + \mathrm{i}\sin \frac{3\pi}{16}\right )\)}{\(\mathrm{i}\root{4}\of{-1-\mathrm{i}}\)}{\(\root{4}\of{-1-\mathrm{i}}\)}{\(\root{8}\of{2}\left (\cos \frac{5\pi}{16} + \mathrm{i}\sin \frac{5\pi}{16}\right )\)}{}{}}
\End

\Data\ID{9000034301}
\Updated{2022-04-23 05:04:23}
\Level{A}
\SubArea{Binomial equations}
\LANGen{
\Question{
Find the solution set of the following equation in the set of complex numbers.
\[ 
 x^{3} - 1 = 0
\]}
{\(\{1;\ -\frac{1}
{2} + \mathrm{i}\frac{\sqrt{3}} 
 {2} ;\ -\frac{1}
{2} -\mathrm{i}\frac{\sqrt{3}} 
 {2} \}\)}{\(\{1;\ -1 + \mathrm{i}\sqrt{3};\ -1 -\mathrm{i}\sqrt{3}\}\)}{\(\{1;\ -\frac{1}
{2} + \mathrm{i}\frac{\sqrt{3}} 
 {2} \}\)}{\(\{1;\ -\frac{1}
{2} -\mathrm{i}\frac{\sqrt{3}} 
 {2} \}\)}{}{}}
\LANGpl{
\Question{
Wyznacz zbiór rozwiązań równania w zbiorze liczb zespolonych.
\[ 
 x^{3} - 1 = 0
\]}
{\(\{1;\ -\frac{1}
{2} + \mathrm{i}\frac{\sqrt{3}} 
 {2} ;\ -\frac{1}
{2} -\mathrm{i}\frac{\sqrt{3}} 
 {2} \}\)}{\(\{1;\ -1 + \mathrm{i}\sqrt{3};\ -1 -\mathrm{i}\sqrt{3}\}\)}{\(\{1;\ -\frac{1}
{2} + \mathrm{i}\frac{\sqrt{3}} 
 {2} \}\)}{\(\{1;\ -\frac{1}
{2} -\mathrm{i}\frac{\sqrt{3}} 
 {2} \}\)}{}{}}
\LANGcs{
\Question{
Množinou všech komplexních řešení rovnice
\(x^{3} - 1 = 0\) je:}
{\(\{1;\ -\frac{1}
{2} + \mathrm{i}\frac{\sqrt{3}} 
 {2} ;\ -\frac{1}
{2} -\mathrm{i}\frac{\sqrt{3}} 
 {2} \}\)}{\(\{1;\ -1 + \mathrm{i}\sqrt{3};\ -1 -\mathrm{i}\sqrt{3}\}\)}{\(\{1;\ -\frac{1}
{2} + \mathrm{i}\frac{\sqrt{3}} 
 {2} \}\)}{\(\{1;\ -\frac{1}
{2} -\mathrm{i}\frac{\sqrt{3}} 
 {2} \}\)}{}{}}
\LANGsk{
\Question{
Nájdite množinu všetkých riešení danej rovnice v množine komplexných čísel.
\[ 
 x^{3} - 1 = 0
\]}
{\(\{1;\ -\frac{1}
{2} + \mathrm{i}\frac{\sqrt{3}} 
 {2} ;\ -\frac{1}
{2} -\mathrm{i}\frac{\sqrt{3}} 
 {2} \}\)}{\(\{1;\ -1 + \mathrm{i}\sqrt{3};\ -1 -\mathrm{i}\sqrt{3}\}\)}{\(\{1;\ -\frac{1}
{2} + \mathrm{i}\frac{\sqrt{3}} 
 {2} \}\)}{\(\{1;\ -\frac{1}
{2} -\mathrm{i}\frac{\sqrt{3}} 
 {2} \}\)}{}{}}
\LANGes{
\Question{
Determina el conjunto de soluciones de la siguiente ecuación en el conjunto de los números complejos.
\[ 
 x^{3} - 1 = 0
\]}
{\(\{1;\ -\frac{1}
{2} + \mathrm{i}\frac{\sqrt{3}} 
 {2} ;\ -\frac{1}
{2} -\mathrm{i}\frac{\sqrt{3}} 
 {2} \}\)}{\(\{1;\ -1 + \mathrm{i}\sqrt{3};\ -1 -\mathrm{i}\sqrt{3}\}\)}{\(\{1;\ -\frac{1}
{2} + \mathrm{i}\frac{\sqrt{3}} 
 {2} \}\)}{\(\{1;\ -\frac{1}
{2} -\mathrm{i}\frac{\sqrt{3}} 
 {2} \}\)}{}{}}
\End

\Data\ID{9000034302}
\Updated{2022-04-23 05:04:23}
\Level{A}
\SubArea{Binomial equations}
\LANGen{
\Question{
Find the solution set of the following equation in the set of complex numbers.
\[ 
 x^{3} + 8 = 0
\]}
{\(\{ - 2;\ 1 + \mathrm{i}\sqrt{3};\ 1 -\mathrm{i}\sqrt{3}\}\)}{\(\{ - 2;\ -1 + \mathrm{i}\sqrt{3};\ -1 -\mathrm{i}\sqrt{3}\}\)}{\(\{ - 2;\ \frac{1} 
{2} + \mathrm{i}\frac{\sqrt{3}} 
 {2} ;\ \frac{1} 
{2} -\mathrm{i}\frac{\sqrt{3}} 
 {2} \}\)}{\(\{ - 2;\ -\frac{1}
{2} + \mathrm{i}\frac{\sqrt{3}} 
 {2} ;\ -\frac{1}
{2} -\mathrm{i}\frac{\sqrt{3}} 
 {2} \}\)}{}{}}
\LANGpl{
\Question{
Wyznacz zbiór rozwiązań równania w zbiorze liczb zespolonych.
\[ 
 x^{3} + 8 = 0
\]}
{\(\{ - 2;\ 1 + \mathrm{i}\sqrt{3};\ 1 -\mathrm{i}\sqrt{3}\}\)}{\(\{ - 2;\ -1 + \mathrm{i}\sqrt{3};\ -1 -\mathrm{i}\sqrt{3}\}\)}{\(\{ - 2;\ \frac{1} 
{2} + \mathrm{i}\frac{\sqrt{3}} 
 {2} ;\ \frac{1} 
{2} -\mathrm{i}\frac{\sqrt{3}} 
 {2} \}\)}{\(\{ - 2;\ -\frac{1}
{2} + \mathrm{i}\frac{\sqrt{3}} 
 {2} ;\ -\frac{1}
{2} -\mathrm{i}\frac{\sqrt{3}} 
 {2} \}\)}{}{}}
\LANGcs{
\Question{
Množinou všech komplexních řešení rovnice
\(x^{3} + 8 = 0\) je:}
{\(\{ - 2;\ 1 + \mathrm{i}\sqrt{3};\ 1 -\mathrm{i}\sqrt{3}\}\)}{\(\{ - 2;\ -1 + \mathrm{i}\sqrt{3};\ -1 -\mathrm{i}\sqrt{3}\}\)}{\(\{ - 2;\ \frac{1} 
{2} + \mathrm{i}\frac{\sqrt{3}} 
 {2} ;\ \frac{1} 
{2} -\mathrm{i}\frac{\sqrt{3}} 
 {2} \}\)}{\(\{ - 2;\ -\frac{1}
{2} + \mathrm{i}\frac{\sqrt{3}} 
 {2} ;\ -\frac{1}
{2} -\mathrm{i}\frac{\sqrt{3}} 
 {2} \}\)}{}{}}
\LANGsk{
\Question{
Nájdite množinu všetkých riešení danej rovnice v množine komplexných čísel.
\[ 
 x^{3} + 8 = 0
\]}
{\(\{ - 2;\ 1 + \mathrm{i}\sqrt{3};\ 1 -\mathrm{i}\sqrt{3}\}\)}{\(\{ - 2;\ -1 + \mathrm{i}\sqrt{3};\ -1 -\mathrm{i}\sqrt{3}\}\)}{\(\{ - 2;\ \frac{1} 
{2} + \mathrm{i}\frac{\sqrt{3}} 
 {2} ;\ \frac{1} 
{2} -\mathrm{i}\frac{\sqrt{3}} 
 {2} \}\)}{\(\{ - 2;\ -\frac{1}
{2} + \mathrm{i}\frac{\sqrt{3}} 
 {2} ;\ -\frac{1}
{2} -\mathrm{i}\frac{\sqrt{3}} 
 {2} \}\)}{}{}}
\LANGes{
\Question{
Determina el conjunto de soluciones de la siguiente ecuación en el conjunto de los números complejos.
\[ 
 x^{3} + 8 = 0
\]}
{\(\{ - 2;\ 1 + \mathrm{i}\sqrt{3};\ 1 -\mathrm{i}\sqrt{3}\}\)}{\(\{ - 2;\ -1 + \mathrm{i}\sqrt{3};\ -1 -\mathrm{i}\sqrt{3}\}\)}{\(\{ - 2;\ \frac{1} 
{2} + \mathrm{i}\frac{\sqrt{3}} 
 {2} ;\ \frac{1} 
{2} -\mathrm{i}\frac{\sqrt{3}} 
 {2} \}\)}{\(\{ - 2;\ -\frac{1}
{2} + \mathrm{i}\frac{\sqrt{3}} 
 {2} ;\ -\frac{1}
{2} -\mathrm{i}\frac{\sqrt{3}} 
 {2} \}\)}{}{}}
\End

\Data\ID{9000034303}
\Updated{2022-04-23 05:04:23}
\Level{A}
\SubArea{Binomial equations}
\LANGen{
\Question{
Find the solution set of the following equation in the set of complex numbers.
\[ 
 x^{3} + \mathrm{i} = 0
\]}
{\(\{\mathrm{i};\ \frac{\sqrt{3}} 
 {2} -\frac{1} 
{2}\mathrm{i};\ -\frac{\sqrt{3}}
 {2} -\frac{1} 
{2}\mathrm{i}\}\)}{\(\{ - 1;\ -\frac{\sqrt{3}}
 {2} + \frac{1} 
{2}\mathrm{i};\ -\frac{\sqrt{3}}
 {2} -\frac{1} 
{2}\mathrm{i}\}\)}{\(\{ - 1;\ \frac{\sqrt{3}} 
 {2} -\frac{1} 
{2}\mathrm{i};\ -\frac{\sqrt{3}}
 {2} -\frac{1} 
{2}\mathrm{i}\}\)}{\(\{\mathrm{i};\ -\frac{\sqrt{3}}
 {2} + \frac{1} 
{2}\mathrm{i};\ -\frac{\sqrt{3}}
 {2} -\frac{1} 
{2}\mathrm{i}\}\)}{}{}}
\LANGpl{
\Question{
Wyznacz zbiór rozwiązań równania w zbiorze liczb zespolonych.
\[ 
 x^{3} + \mathrm{i} = 0
\]}
{\(\{\mathrm{i};\ \frac{\sqrt{3}} 
 {2} -\frac{1} 
{2}\mathrm{i};\ -\frac{\sqrt{3}}
 {2} -\frac{1} 
{2}\mathrm{i}\}\)}{\(\{ - 1;\ -\frac{\sqrt{3}}
 {2} + \frac{1} 
{2}\mathrm{i};\ -\frac{\sqrt{3}}
 {2} -\frac{1} 
{2}\mathrm{i}\}\)}{\(\{ - 1;\ \frac{\sqrt{3}} 
 {2} -\frac{1} 
{2}\mathrm{i};\ -\frac{\sqrt{3}}
 {2} -\frac{1} 
{2}\mathrm{i}\}\)}{\(\{\mathrm{i};\ -\frac{\sqrt{3}}
 {2} + \frac{1} 
{2}\mathrm{i};\ -\frac{\sqrt{3}}
 {2} -\frac{1} 
{2}\mathrm{i}\}\)}{}{}}
\LANGcs{
\Question{
Množinou všech komplexních řešení rovnice
\(x^{3} + \mathrm{i} = 0\) je:}
{\(\{\mathrm{i};\ \frac{\sqrt{3}} 
 {2} -\frac{1} 
{2}\mathrm{i};\ -\frac{\sqrt{3}}
 {2} -\frac{1} 
{2}\mathrm{i}\}\)}{\(\{ - 1;\ -\frac{\sqrt{3}}
 {2} + \frac{1} 
{2}\mathrm{i};\ -\frac{\sqrt{3}}
 {2} -\frac{1} 
{2}\mathrm{i}\}\)}{\(\{ - 1;\ \frac{\sqrt{3}} 
 {2} -\frac{1} 
{2}\mathrm{i};\ -\frac{\sqrt{3}}
 {2} -\frac{1} 
{2}\mathrm{i}\}\)}{\(\{\mathrm{i};\ -\frac{\sqrt{3}}
 {2} + \frac{1} 
{2}\mathrm{i};\ -\frac{\sqrt{3}}
 {2} -\frac{1} 
{2}\mathrm{i}\}\)}{}{}}
\LANGsk{
\Question{
Nájdite množinu všetkých riešení danej rovnice v množine komplexných čísel. 
\[ 
 x^{3} + \mathrm{i} = 0
\]}
{\(\{\mathrm{i};\ \frac{\sqrt{3}} 
 {2} -\frac{1} 
{2}\mathrm{i};\ -\frac{\sqrt{3}}
 {2} -\frac{1} 
{2}\mathrm{i}\}\)}{\(\{ - 1;\ -\frac{\sqrt{3}}
 {2} + \frac{1} 
{2}\mathrm{i};\ -\frac{\sqrt{3}}
 {2} -\frac{1} 
{2}\mathrm{i}\}\)}{\(\{ - 1;\ \frac{\sqrt{3}} 
 {2} -\frac{1} 
{2}\mathrm{i};\ -\frac{\sqrt{3}}
 {2} -\frac{1} 
{2}\mathrm{i}\}\)}{\(\{\mathrm{i};\ -\frac{\sqrt{3}}
 {2} + \frac{1} 
{2}\mathrm{i};\ -\frac{\sqrt{3}}
 {2} -\frac{1} 
{2}\mathrm{i}\}\)}{}{}}
\LANGes{
\Question{
Determina el conjunto de soluciones de la siguiente ecuación en el conjunto de los números complejos.
\[ 
 x^{3} + \mathrm{i} = 0
\]}
{\(\{\mathrm{i};\ \frac{\sqrt{3}} 
 {2} -\frac{1} 
{2}\mathrm{i};\ -\frac{\sqrt{3}}
 {2} -\frac{1} 
{2}\mathrm{i}\}\)}{\(\{ - 1;\ -\frac{\sqrt{3}}
 {2} + \frac{1} 
{2}\mathrm{i};\ -\frac{\sqrt{3}}
 {2} -\frac{1} 
{2}\mathrm{i}\}\)}{\(\{ - 1;\ \frac{\sqrt{3}} 
 {2} -\frac{1} 
{2}\mathrm{i};\ -\frac{\sqrt{3}}
 {2} -\frac{1} 
{2}\mathrm{i}\}\)}{\(\{\mathrm{i};\ -\frac{\sqrt{3}}
 {2} + \frac{1} 
{2}\mathrm{i};\ -\frac{\sqrt{3}}
 {2} -\frac{1} 
{2}\mathrm{i}\}\)}{}{}}
\End

\Data\ID{9000034304}
\Updated{2021-04-25 06:04:42}
\Level{A}
\SubArea{Binomial equations}
\LANGen{
\Question{
Find the solution set of the following equation in the set of complex numbers.
\[ 
 x^{4} - 1 = 0
\]}
{\(\{1;\ -1;\ \mathrm{i};\ -\mathrm{i}\}\)}{\(\{1 -\mathrm{i};\ -1 -\mathrm{i}\}\)}{\(\{1 + \mathrm{i};\ -1 + \mathrm{i}\}\)}{\(\{1 + \mathrm{i};\ 1 -\mathrm{i};\ -1 + \mathrm{i};\ -1 -\mathrm{i}\}\)}{}{}}
\LANGpl{
\Question{
Wyznacz zbiór rozwiązań równania w zbiorze liczb zespolonych.
\[ 
 x^{4} - 1 = 0
\]}
{\(\{1;\ -1;\ \mathrm{i};\ -\mathrm{i}\}\)}{\(\{1 -\mathrm{i};\ -1 -\mathrm{i}\}\)}{\(\{1 + \mathrm{i};\ -1 + \mathrm{i}\}\)}{\(\{1 + \mathrm{i};\ 1 -\mathrm{i};\ -1 + \mathrm{i};\ -1 -\mathrm{i}\}\)}{}{}}
\LANGcs{
\Question{
Množinou všech komplexních řešení rovnice
\(x^{4} - 1 = 0\) je:}
{\(\{1;\ -1;\ \mathrm{i};\ -\mathrm{i}\}\)}{\(\{1 -\mathrm{i};\ -1 -\mathrm{i}\}\)}{\(\{1 + \mathrm{i};\ -1 + \mathrm{i}\}\)}{\(\{1 + \mathrm{i};\ 1 -\mathrm{i};\ -1 + \mathrm{i};\ -1 -\mathrm{i}\}\)}{}{}}
\LANGsk{
\Question{
Nájdite množinu všetkých riešení danej rovnice v množine komplexných čísel. 
\[ 
 x^{4} - 1 = 0
\]}
{\(\{1;\ -1;\ \mathrm{i};\ -\mathrm{i}\}\)}{\(\{1 -\mathrm{i};\ -1 -\mathrm{i}\}\)}{\(\{1 + \mathrm{i};\ -1 + \mathrm{i}\}\)}{\(\{1 + \mathrm{i};\ 1 -\mathrm{i};\ -1 + \mathrm{i};\ -1 -\mathrm{i}\}\)}{}{}}
\LANGes{
\Question{
Determina el conjunto de soluciones de la siguiente ecuación en el conjunto de los números complejos.
\[ 
 x^{4} - 1 = 0
\]}
{\(\{1;\ -1;\ \mathrm{i};\ -\mathrm{i}\}\)}{\(\{1 -\mathrm{i};\ -1 -\mathrm{i}\}\)}{\(\{1 + \mathrm{i};\ -1 + \mathrm{i}\}\)}{\(\{1 + \mathrm{i};\ 1 -\mathrm{i};\ -1 + \mathrm{i};\ -1 -\mathrm{i}\}\)}{}{}}
\End

\Data\ID{9000034305}
\Updated{2021-04-25 06:04:20}
\Level{A}
\SubArea{Binomial equations}
\LANGen{
\Question{
Solve the following equation in the set of complex numbers. 
\[ 
 x^{4} + 16 = 0
\]}
{\(x_{1, 2} = \sqrt{2}(1\pm \mathrm{i}),\ x_{3, 4} = -\sqrt{2}(1\pm \mathrm{i})\)}{\(x_{1, 2} = 1\pm \mathrm{i},\ x_{3, 4} = -1\pm \mathrm{i}\)}{\(x_{1, 2} = 2(1\pm \mathrm{i}),\ x_{3, 4} = -2(1\pm \mathrm{i})\)}{\(x_{1, 2} = \frac{\sqrt{2}} 
 {2} (1\pm \mathrm{i}),\ x_{3, 4} = -\frac{\sqrt{2}}
 {2} (1\pm \mathrm{i})\)}{}{}}
\LANGpl{
\Question{
Rozwiąż równanie w zbiorze liczb zespolonych.
\[ 
 x^{4} + 16 = 0
\]}
{\(x_{1, 2} = \sqrt{2}(1\pm \mathrm{i}),\ x_{3, 4} = -\sqrt{2}(1\pm \mathrm{i})\)}{\(x_{1, 2} = 1\pm \mathrm{i},\ x_{3, 4} = -1\pm \mathrm{i}\)}{\(x_{1, 2} = 2(1\pm \mathrm{i}),\ x_{3, 4} = -2(1\pm \mathrm{i})\)}{\(x_{1, 2} = \frac{\sqrt{2}} 
 {2} (1\pm \mathrm{i}),\ x_{3, 4} = -\frac{\sqrt{2}}
 {2} (1\pm \mathrm{i})\)}{}{}}
\LANGcs{
\Question{
Která z následujících možností vyjadřuje všechna řešení rovnice
\(x^{4} + 16 = 0\) s neznámou
\(x\in \mathbb{C}\)?}
{\(x_{1, 2} = \sqrt{2}(1\pm \mathrm{i}),\ x_{3, 4} = -\sqrt{2}(1\pm \mathrm{i})\)}{\(x_{1, 2} = 1\pm \mathrm{i},\ x_{3, 4} = -1\pm \mathrm{i}\)}{\(x_{1, 2} = 2(1\pm \mathrm{i}),\ x_{3, 4} = -2(1\pm \mathrm{i})\)}{\(x_{1, 2} = \frac{\sqrt{2}} 
 {2} (1\pm \mathrm{i}),\ x_{3, 4} = -\frac{\sqrt{2}}
 {2} (1\pm \mathrm{i})\)}{}{}}
\LANGsk{
\Question{
Vyriešte danú rovnicu v množine komplexných čísel. 
\[ 
 x^{4} + 16 = 0
\]}
{\(x_{1, 2} = \sqrt{2}(1\pm \mathrm{i}),\ x_{3, 4} = -\sqrt{2}(1\pm \mathrm{i})\)}{\(x_{1, 2} = 1\pm \mathrm{i},\ x_{3, 4} = -1\pm \mathrm{i}\)}{\(x_{1, 2} = 2(1\pm \mathrm{i}),\ x_{3, 4} = -2(1\pm \mathrm{i})\)}{\(x_{1, 2} = \frac{\sqrt{2}} 
 {2} (1\pm \mathrm{i}),\ x_{3, 4} = -\frac{\sqrt{2}}
 {2} (1\pm \mathrm{i})\)}{}{}}
\LANGes{
\Question{
Resuelve la siguiente ecuación en el conjunto de los números complejos.
\[ 
 x^{4} + 16 = 0
\]}
{\(x_{1, 2} = \sqrt{2}(1\pm \mathrm{i}),\ x_{3, 4} = -\sqrt{2}(1\pm \mathrm{i})\)}{\(x_{1, 2} = 1\pm \mathrm{i},\ x_{3, 4} = -1\pm \mathrm{i}\)}{\(x_{1, 2} = 2(1\pm \mathrm{i}),\ x_{3, 4} = -2(1\pm \mathrm{i})\)}{\(x_{1, 2} = \frac{\sqrt{2}} 
 {2} (1\pm \mathrm{i}),\ x_{3, 4} = -\frac{\sqrt{2}}
 {2} (1\pm \mathrm{i})\)}{}{}}
\End

\Data\ID{9000034306}
\Updated{2021-04-25 06:04:40}
\Level{A}
\SubArea{Binomial equations}
\LANGen{
\Question{
Solve the following equation in the set of complex numbers. 
\[ 
 x^{6} - 64 = 0
\]}
{\(x_{1, 2} =\pm 2,\ x_{3, 4} = 1\pm \mathrm{i}\sqrt{3},\ x_{5, 6} = -1\pm \mathrm{i}\sqrt{3}\)}{\(x_{1, 2} =\pm 2,\ x_{3, 4} = \frac{1} 
{2}\pm \mathrm{i}\frac{\sqrt{3}} 
 {2} ,\ x_{5, 6} = -\frac{1}
{2}\pm \mathrm{i}\frac{\sqrt{3}} 
 {2} \)}{\(x_{1, 2} =\pm 4,\ x_{3, 4} = 1\pm \mathrm{i}\sqrt{3},\ x_{5, 6} = -1\pm \mathrm{i}\sqrt{3}\)}{\(x_{1, 2} =\pm 8,\ x_{3, 4} = 2\pm 2\mathrm{i}\sqrt{3},\ x_{5, 6} = -2\pm 2\mathrm{i}\sqrt{3}\)}{}{}}
\LANGpl{
\Question{
Rozwiąż równanie w zbiorze liczb zespolonych.
\[ 
 x^{6} - 64 = 0
\]}
{\(x_{1, 2} =\pm 2,\ x_{3, 4} = 1\pm \mathrm{i}\sqrt{3},\ x_{5, 6} = -1\pm \mathrm{i}\sqrt{3}\)}{\(x_{1, 2} =\pm 2,\ x_{3, 4} = \frac{1} 
{2}\pm \mathrm{i}\frac{\sqrt{3}} 
 {2} ,\ x_{5, 6} = -\frac{1}
{2}\pm \mathrm{i}\frac{\sqrt{3}} 
 {2} \)}{\(x_{1, 2} =\pm 4,\ x_{3, 4} = 1\pm \mathrm{i}\sqrt{3},\ x_{5, 6} = -1\pm \mathrm{i}\sqrt{3}\)}{\(x_{1, 2} =\pm 8,\ x_{3, 4} = 2\pm 2\mathrm{i}\sqrt{3},\ x_{5, 6} = -2\pm 2\mathrm{i}\sqrt{3}\)}{}{}}
\LANGcs{
\Question{
Která z následujících možností vyjadřuje všechna řešení rovnice
\(x^{6} - 64 = 0\) s neznámou
\(x\in \mathbb{C}\)?}
{\(x_{1, 2} =\pm 2,\ x_{3, 4} = 1\pm \mathrm{i}\sqrt{3},\ x_{5, 6} = -1\pm \mathrm{i}\sqrt{3}\)}{\(x_{1, 2} =\pm 2,\ x_{3, 4} = \frac{1} 
{2}\pm \mathrm{i}\frac{\sqrt{3}} 
 {2} ,\ x_{5, 6} = -\frac{1}
{2}\pm \mathrm{i}\frac{\sqrt{3}} 
 {2} \)}{\(x_{1, 2} =\pm 4,\ x_{3, 4} = 1\pm \mathrm{i}\sqrt{3},\ x_{5, 6} = -1\pm \mathrm{i}\sqrt{3}\)}{\(x_{1, 2} =\pm 8,\ x_{3, 4} = 2\pm 2\mathrm{i}\sqrt{3},\ x_{5, 6} = -2\pm 2\mathrm{i}\sqrt{3}\)}{}{}}
\LANGsk{
\Question{
Vyriešte nasledujúcu rovnicu v množine komplexných čísel. 
\[ 
 x^{6} - 64 = 0
\]}
{\(x_{1, 2} =\pm 2,\ x_{3, 4} = 1\pm \mathrm{i}\sqrt{3},\ x_{5, 6} = -1\pm \mathrm{i}\sqrt{3}\)}{\(x_{1, 2} =\pm 2,\ x_{3, 4} = \frac{1} 
{2}\pm \mathrm{i}\frac{\sqrt{3}} 
 {2} ,\ x_{5, 6} = -\frac{1}
{2}\pm \mathrm{i}\frac{\sqrt{3}} 
 {2} \)}{\(x_{1, 2} =\pm 4,\ x_{3, 4} = 1\pm \mathrm{i}\sqrt{3},\ x_{5, 6} = -1\pm \mathrm{i}\sqrt{3}\)}{\(x_{1, 2} =\pm 8,\ x_{3, 4} = 2\pm 2\mathrm{i}\sqrt{3},\ x_{5, 6} = -2\pm 2\mathrm{i}\sqrt{3}\)}{}{}}
\LANGes{
\Question{
Resuelve la siguiente ecuación en el conjunto de los números complejos.
\[ 
 x^{6} - 64 = 0
\]}
{\(x_{1, 2} =\pm 2,\ x_{3, 4} = 1\pm \mathrm{i}\sqrt{3},\ x_{5, 6} = -1\pm \mathrm{i}\sqrt{3}\)}{\(x_{1, 2} =\pm 2,\ x_{3, 4} = \frac{1} 
{2}\pm \mathrm{i}\frac{\sqrt{3}} 
 {2} ,\ x_{5, 6} = -\frac{1}
{2}\pm \mathrm{i}\frac{\sqrt{3}} 
 {2} \)}{\(x_{1, 2} =\pm 4,\ x_{3, 4} = 1\pm \mathrm{i}\sqrt{3},\ x_{5, 6} = -1\pm \mathrm{i}\sqrt{3}\)}{\(x_{1, 2} =\pm 8,\ x_{3, 4} = 2\pm 2\mathrm{i}\sqrt{3},\ x_{5, 6} = -2\pm 2\mathrm{i}\sqrt{3}\)}{}{}}
\End

\Data\ID{9000034307}
\Updated{2020-07-19 02:07:56}
\Level{A}
\SubArea{Binomial equations}
\LANGen{
\Question{
Let \(x\) 
be any of the complex solutions of the following equation. 
\[ 
 x^{5} + \sqrt{3} -\mathrm{i} = 0
\]
Find the absolute value of \(x\).}
{\(\root{5}\of{2}\)}{\(2\)}{\(\root{5}\of{4}\)}{\(\root{10}\of{2}\)}{}{}}
\LANGpl{
\Question{
Wartość bezwzględna każdego rozwiązania równania 
\[ 
 x^{5} + \sqrt{3} -\mathrm{i} = 0
\]
wynosi:}
{\(\root{5}\of{2}\)}{\(2\)}{\(\root{5}\of{4}\)}{\(\root{10}\of{2}\)}{}{}}
\LANGcs{
\Question{
Absolutní hodnota každého řešení rovnice
\(x^{5} + \sqrt{3} -\mathrm{i} = 0\) je:}
{\(\root{5}\of{2}\)}{\(2\)}{\(\root{5}\of{4}\)}{\(\root{10}\of{2}\)}{}{}}
\LANGsk{
\Question{
Nech \(x\) 
je jedno z komplexných riešení danej rovnice.  
\[ 
 x^{5} + \sqrt{3} -\mathrm{i} = 0
\]
Nájdite absolútnu hodnotu \(x\).}
{\(\root{5}\of{2}\)}{\(2\)}{\(\root{5}\of{4}\)}{\(\root{10}\of{2}\)}{}{}}
\LANGes{
\Question{
Sea \(x\) 
cualquiera de las soluciones complejas de la siguiente ecuación.
\[ 
 x^{5} + \sqrt{3} -\mathrm{i} = 0
\]
Calcula el valor absoluto de \(x\).}
{\(\root{5}\of{2}\)}{\(2\)}{\(\root{5}\of{4}\)}{\(\root{10}\of{2}\)}{}{}}
\End

\Data\ID{9000034308}
\Updated{2022-04-23 05:04:23}
\Level{A}
\SubArea{Binomial equations}
\LANGen{
\Question{
Two solutions of the equation 
\[ 
 x^{3} + 1 + \mathrm{i} = 0
\]
are 
\[ 
\begin{aligned}x_{1}& = \root{6}\of{2}\left (\cos \frac{5}
{12}\pi + \mathrm{i}\sin \frac{5}
{12}\pi \right ),&
\\x_{2}& = \root{6}\of{2}\left (\cos \frac{13}
{12}\pi + \mathrm{i}\sin \frac{13}
{12}\pi \right ).
\\ \end{aligned}
\]
Find the third solution.}
{\(x_{3} = \root{6}\of{2}\left (\cos \frac{21}
{12}\pi + \mathrm{i}\sin \frac{21}
{12}\pi \right )\)}{\(x_{3} = \root{6}\of{2}\left (\cos \frac{9}
{12}\pi + \mathrm{i}\sin \frac{9}
{12}\pi \right )\)}{\(x_{3} = \root{6}\of{2}\left (\cos \frac{17}
{12}\pi + \mathrm{i}\sin \frac{17}
{12}\pi \right )\)}{\(x_{3} = \root{6}\of{2}\left (\cos \frac{19}
{12}\pi + \mathrm{i}\sin \frac{19}
{12}\pi \right )\)}{}{}}
\LANGpl{
\Question{
Równanie
\[ 
 x^{3} + 1 + \mathrm{i} = 0
\]
ma dwa rozwiązania
\[ 
\begin{aligned}x_{1}& = \root{6}\of{2}\left (\cos \frac{5}
{12}\pi + \mathrm{i}\sin \frac{5}
{12}\pi \right ),&
\\x_{2}& = \root{6}\of{2}\left (\cos \frac{13}
{12}\pi + \mathrm{i}\sin \frac{13}
{12}\pi \right ).
\\ \end{aligned}
\]
Wskaż trzecie.}
{\(x_{3} = \root{6}\of{2}\left (\cos \frac{21}
{12}\pi + \mathrm{i}\sin \frac{21}
{12}\pi \right )\)}{\(x_{3} = \root{6}\of{2}\left (\cos \frac{9}
{12}\pi + \mathrm{i}\sin \frac{9}
{12}\pi \right )\)}{\(x_{3} = \root{6}\of{2}\left (\cos \frac{17}
{12}\pi + \mathrm{i}\sin \frac{17}
{12}\pi \right )\)}{\(x_{3} = \root{6}\of{2}\left (\cos \frac{19}
{12}\pi + \mathrm{i}\sin \frac{19}
{12}\pi \right )\)}{}{}}
\LANGcs{
\Question{
Dvě z řešení rovnice
 \[x^{3} + 1 + \mathrm{i} = 0\]
jsou
\[ 
 x_{1} = \root{6}\of{2}\left (\cos \frac{5}
{12}\pi + \mathrm{i}\sin \frac{5}
{12}\pi \right ),
\]
\[ 
 x_{2} = \root{6}\of{2}\left (\cos \frac{13}
{12}\pi + \mathrm{i}\sin \frac{13}
{12}\pi \right ).
\]
 Třetím řešení rovnice je:}
{\(x_{3} = \root{6}\of{2}\left (\cos \frac{21}
{12}\pi + \mathrm{i}\sin \frac{21}
{12}\pi \right )\)}{\(x_{3} = \root{6}\of{2}\left (\cos \frac{9}
{12}\pi + \mathrm{i}\sin \frac{9}
{12}\pi \right )\)}{\(x_{3} = \root{6}\of{2}\left (\cos \frac{17}
{12}\pi + \mathrm{i}\sin \frac{17}
{12}\pi \right )\)}{\(x_{3} = \root{6}\of{2}\left (\cos \frac{19}
{12}\pi + \mathrm{i}\sin \frac{19}
{12}\pi \right )\)}{}{}}
\LANGsk{
\Question{
Dve riešenia rovnice
\[ 
 x^{3} + 1 + \mathrm{i} = 0
\]
sú 
\[ 
\begin{aligned}x_{1}& = \root{6}\of{2}\left (\cos \frac{5}
{12}\pi + \mathrm{i}\sin \frac{5}
{12}\pi \right ),&
\\x_{2}& = \root{6}\of{2}\left (\cos \frac{13}
{12}\pi + \mathrm{i}\sin \frac{13}
{12}\pi \right ).
\\ \end{aligned}
\]
Nájdite tretie riešenie.}
{\(x_{3} = \root{6}\of{2}\left (\cos \frac{21}
{12}\pi + \mathrm{i}\sin \frac{21}
{12}\pi \right )\)}{\(x_{3} = \root{6}\of{2}\left (\cos \frac{9}
{12}\pi + \mathrm{i}\sin \frac{9}
{12}\pi \right )\)}{\(x_{3} = \root{6}\of{2}\left (\cos \frac{17}
{12}\pi + \mathrm{i}\sin \frac{17}
{12}\pi \right )\)}{\(x_{3} = \root{6}\of{2}\left (\cos \frac{19}
{12}\pi + \mathrm{i}\sin \frac{19}
{12}\pi \right )\)}{}{}}
\LANGes{
\Question{
Dos soluciones de la ecuación
\[ 
 x^{3} + 1 + \mathrm{i} = 0
\]
son
\[ 
\begin{aligned}x_{1}& = \root{6}\of{2}\left (\cos \frac{5}
{12}\pi + \mathrm{i}\sin \frac{5}
{12}\pi \right ),&
\\x_{2}& = \root{6}\of{2}\left (\cos \frac{13}
{12}\pi + \mathrm{i}\sin \frac{13}
{12}\pi \right ).
\\ \end{aligned}
\]
Calcula la tercera solución.}
{\(x_{3} = \root{6}\of{2}\left (\cos \frac{21}
{12}\pi + \mathrm{i}\sin \frac{21}
{12}\pi \right )\)}{\(x_{3} = \root{6}\of{2}\left (\cos \frac{9}
{12}\pi + \mathrm{i}\sin \frac{9}
{12}\pi \right )\)}{\(x_{3} = \root{6}\of{2}\left (\cos \frac{17}
{12}\pi + \mathrm{i}\sin \frac{17}
{12}\pi \right )\)}{\(x_{3} = \root{6}\of{2}\left (\cos \frac{19}
{12}\pi + \mathrm{i}\sin \frac{19}
{12}\pi \right )\)}{}{}}
\End

\Data\ID{9000034309}
\Updated{2021-04-25 06:04:22}
\Level{A}
\SubArea{Binomial equations}
\LANGen{
\Question{
Find the angle \(\varphi \) 
such that the angles in the polar form of any two solutions of the equation
\[ 
 x^{5} - 1 + \mathrm{i}\sqrt{3} = 0
\]
differ by an integer multiple of \(\varphi \).}
{\(\varphi = \frac{2} 
{5}\pi \)}{\(\varphi = \frac{3} 
{5}\pi \)}{\(\varphi = \frac{4} 
{5}\pi \)}{\(\varphi =\pi \)}{}{}}
\LANGpl{
\Question{
Wskaż kąt  \(\varphi \) 
tak, aby kąty  w postaci biegunowej jakichkolwiek dwóch rozwiązań równania
\[ 
 x^{5} - 1 + \mathrm{i}\sqrt{3} = 0
\]
różniły się liczbą całkowitą wielokrotności  \(\varphi \).}
{\(\varphi = \frac{2} 
{5}\pi \)}{\(\varphi = \frac{3} 
{5}\pi \)}{\(\varphi = \frac{4} 
{5}\pi \)}{\(\varphi =\pi \)}{}{}}
\LANGcs{
\Question{
Argumenty libovolných dvou řešení rovnice
\(x^{5} - 1 + \mathrm{i}\sqrt{3} = 0\) se
liší o celočíselný násobek čísla:}
{\(\varphi=\frac{2}
{5}\pi \)}{\(\varphi=\frac{3}
{5}\pi \)}{\(\varphi=\frac{4}
{5}\pi \)}{\(\varphi=\pi \)}{}{}}
\LANGsk{
\Question{
Argumenty ľubovoľných dvoch riešení rovnice 
\[ 
 x^{5} - 1 + \mathrm{i}\sqrt{3} = 0
\]
sa líšia o celočíselný násobok čísla:}
{\(\varphi = \frac{2} 
{5}\pi \)}{\(\varphi = \frac{3} 
{5}\pi \)}{\(\varphi = \frac{4} 
{5}\pi \)}{\(\varphi =\pi \)}{}{}}
\LANGes{
\Question{
Calcula el ángulo \(\varphi \) 
suponiendo que los ángulos en la forma polar de cualquiera de las dos soluciones de la ecuación
\[ 
 x^{5} - 1 + \mathrm{i}\sqrt{3} = 0
\]
difieren por un múltiplo entero de \(\varphi \).}
{\(\varphi = \frac{2} 
{5}\pi \)}{\(\varphi = \frac{3} 
{5}\pi \)}{\(\varphi = \frac{4} 
{5}\pi \)}{\(\varphi =\pi \)}{}{}}
\End

\Data\ID{9000035810}
\Updated{2020-07-19 02:07:36}
\Level{A}
\SubArea{Binomial equations}
\LANGen{
\Question{
Given the complex number \(z = -2 + 2\mathrm{i}\),
find all the roots of \(\root{3}\of{z}\).}
{\(\begin{aligned}[t] &w_{0} = \root{6}\of{8}\left (\cos \frac{\pi }
{4} + \mathrm{i}\sin \frac{\pi }
{4}\right ) &
\\&w_{1} = \root{6}\of{8}\left (\cos \frac{11\pi }
{12} + \mathrm{i}\sin \frac{11\pi }
{12}\right )
\\&w_{2} = \root{6}\of{8}\left (\cos \frac{19\pi }
{12} + \mathrm{i}\sin \frac{19\pi }
{12}\right )
\\ \end{aligned}\)}{\(\begin{aligned}[t] &w_{0} = 2\left (\cos \frac{\pi }
{4} + \mathrm{i}\sin \frac{\pi }
{4}\right ) &
\\&w_{1} = 2\left (\cos \frac{11\pi }
{12} + \mathrm{i}\sin \frac{11\pi }
{12}\right )
\\&w_{2} = 2\left (\cos \frac{19\pi }
{12} + \mathrm{i}\sin \frac{19\pi }
{12}\right )
\\ \end{aligned}\)}{\(\root{3}\of{-2} + \root{3}\of{2}\)}{\(\begin{aligned}[t] &w_{0} = 2\left (\cos \frac{\pi }
{3} + \mathrm{i}\sin \frac{\pi }
{3}\right )&
\\&w_{1} = 2\left (\cos \pi +\mathrm{i}\sin \pi \right )
\\&w_{2} = 2\left (\cos \frac{5\pi }
{3} + \mathrm{i}\sin \frac{5\pi }
{3}\right )
\\ \end{aligned}\)}{}{}}
\LANGpl{
\Question{
Podano liczbę zespoloną \(z = -2 + 2\mathrm{i}\),
wskaż wszystkie pierwiastki \(\root{3}\of{z}\).}
{\(\begin{aligned}[t] &w_{0} = \root{6}\of{8}\left (\cos \frac{\pi }
{4} + \mathrm{i}\sin \frac{\pi }
{4}\right ) &
\\&w_{1} = \root{6}\of{8}\left (\cos \frac{11\pi }
{12} + \mathrm{i}\sin \frac{11\pi }
{12}\right )
\\&w_{2} = \root{6}\of{8}\left (\cos \frac{19\pi }
{12} + \mathrm{i}\sin \frac{19\pi }
{12}\right )
\\ \end{aligned}\)}{\(\begin{aligned}[t] &w_{0} = 2\left (\cos \frac{\pi }
{4} + \mathrm{i}\sin \frac{\pi }
{4}\right ) &
\\&w_{1} = 2\left (\cos \frac{11\pi }
{12} + \mathrm{i}\sin \frac{11\pi }
{12}\right )
\\&w_{2} = 2\left (\cos \frac{19\pi }
{12} + \mathrm{i}\sin \frac{19\pi }
{12}\right )
\\ \end{aligned}\)}{\(\root{3}\of{-2} + \root{3}\of{2}\)}{\(\begin{aligned}[t] &w_{0} = 2\left (\cos \frac{\pi }
{3} + \mathrm{i}\sin \frac{\pi }
{3}\right )&
\\&w_{1} = 2\left (\cos \pi +\mathrm{i}\sin \pi \right )
\\&w_{2} = 2\left (\cos \frac{5\pi }
{3} + \mathrm{i}\sin \frac{5\pi }
{3}\right )
\\ \end{aligned}\)}{}{}}
\LANGcs{
\Question{
Je dáno komplexní číslo \(z = -2 + 2\mathrm{i}\).
Všechny navzájem různé hodnoty \(\root{3}\of{z}\) 
jsou:}
{\(\begin{aligned}[t] &w_{0} = \root{6}\of{8}\left (\cos \frac{\pi }
{4} + \mathrm{i}\sin \frac{\pi }
{4}\right ) &
\\&w_{1} = \root{6}\of{8}\left (\cos \frac{11\pi }
{12} + \mathrm{i}\sin \frac{11\pi }
{12}\right )
\\&w_{2} = \root{6}\of{8}\left (\cos \frac{19\pi }
{12} + \mathrm{i}\sin \frac{19\pi }
{12}\right )
\\ \end{aligned}\)}{\(\begin{aligned}[t] &w_{0} = 2\left (\cos \frac{\pi }
{4} + \mathrm{i}\sin \frac{\pi }
{4}\right ) &
\\&w_{1} = 2\left (\cos \frac{11\pi }
{12} + \mathrm{i}\sin \frac{11\pi }
{12}\right )
\\&w_{2} = 2\left (\cos \frac{19\pi }
{12} + \mathrm{i}\sin \frac{19\pi }
{12}\right )
\\ \end{aligned}\)}{\(\root{3}\of{-2} + \root{3}\of{2}\)}{\(\begin{aligned}[t] &w_{0} = 2\left (\cos \frac{\pi }
{3} + \mathrm{i}\sin \frac{\pi }
{3}\right )&
\\&w_{1} = 2\left (\cos \pi +\mathrm{i}\sin \pi \right )
\\&w_{2} = 2\left (\cos \frac{5\pi }
{3} + \mathrm{i}\sin \frac{5\pi }
{3}\right )
\\ \end{aligned}\)}{}{}}
\LANGsk{
\Question{
Je dané komplexné číslo \(z = -2 + 2\mathrm{i}\).
Všetky navzájom rôzne hodnoty \(\root{3}\of{z}\) sú:}
{\(\begin{aligned}[t] &w_{0} = \root{6}\of{8}\left (\cos \frac{\pi }
{4} + \mathrm{i}\sin \frac{\pi }
{4}\right ) &
\\&w_{1} = \root{6}\of{8}\left (\cos \frac{11\pi }
{12} + \mathrm{i}\sin \frac{11\pi }
{12}\right )
\\&w_{2} = \root{6}\of{8}\left (\cos \frac{19\pi }
{12} + \mathrm{i}\sin \frac{19\pi }
{12}\right )
\\ \end{aligned}\)}{\(\begin{aligned}[t] &w_{0} = 2\left (\cos \frac{\pi }
{4} + \mathrm{i}\sin \frac{\pi }
{4}\right ) &
\\&w_{1} = 2\left (\cos \frac{11\pi }
{12} + \mathrm{i}\sin \frac{11\pi }
{12}\right )
\\&w_{2} = 2\left (\cos \frac{19\pi }
{12} + \mathrm{i}\sin \frac{19\pi }
{12}\right )
\\ \end{aligned}\)}{\(\root{3}\of{-2} + \root{3}\of{2}\)}{\(\begin{aligned}[t] &w_{0} = 2\left (\cos \frac{\pi }
{3} + \mathrm{i}\sin \frac{\pi }
{3}\right )&
\\&w_{1} = 2\left (\cos \pi +\mathrm{i}\sin \pi \right )
\\&w_{2} = 2\left (\cos \frac{5\pi }
{3} + \mathrm{i}\sin \frac{5\pi }
{3}\right )
\\ \end{aligned}\)}{}{}}
\LANGes{
\Question{
Dados los números complejos \(z = -2 + 2\mathrm{i}\),
determina todas las raíces de \(\root{3}\of{z}\).}
{\(\begin{aligned}[t] &w_{0} = \root{6}\of{8}\left (\cos \frac{\pi }
{4} + \mathrm{i}\sin \frac{\pi }
{4}\right ) &
\\&w_{1} = \root{6}\of{8}\left (\cos \frac{11\pi }
{12} + \mathrm{i}\sin \frac{11\pi }
{12}\right )
\\&w_{2} = \root{6}\of{8}\left (\cos \frac{19\pi }
{12} + \mathrm{i}\sin \frac{19\pi }
{12}\right )
\\ \end{aligned}\)}{\(\begin{aligned}[t] &w_{0} = 2\left (\cos \frac{\pi }
{4} + \mathrm{i}\sin \frac{\pi }
{4}\right ) &
\\&w_{1} = 2\left (\cos \frac{11\pi }
{12} + \mathrm{i}\sin \frac{11\pi }
{12}\right )
\\&w_{2} = 2\left (\cos \frac{19\pi }
{12} + \mathrm{i}\sin \frac{19\pi }
{12}\right )
\\ \end{aligned}\)}{\(\root{3}\of{-2} + \root{3}\of{2}\)}{\(\begin{aligned}[t] &w_{0} = 2\left (\cos \frac{\pi }
{3} + \mathrm{i}\sin \frac{\pi }
{3}\right )&
\\&w_{1} = 2\left (\cos \pi +\mathrm{i}\sin \pi \right )
\\&w_{2} = 2\left (\cos \frac{5\pi }
{3} + \mathrm{i}\sin \frac{5\pi }
{3}\right )
\\ \end{aligned}\)}{}{}}
\End
